\documentclass[hidelinks]{article}

% Page Layout
\usepackage[a4paper, total={7in, 10in}]{geometry}

% Colors
\usepackage[dvipsnames,table]{xcolor}

% Math
\usepackage{amsmath,amssymb,mathrsfs}

% Lists
\usepackage{enumitem}

% Tables
\usepackage{array}
\usepackage{tabularx}

% Code listings
\usepackage{listings}

% Section Styling
\usepackage{titlesec}

% Hyperlinks
\usepackage[unicode]{hyperref}
\usepackage[all]{hypcap}

% --------------------------------------------------
% Color Palette
% --------------------------------------------------
\definecolor{carminered}{rgb}{1.0,0.0,0.22}
\definecolor{capri}{rgb}{0.0,0.75,1.0}
\definecolor{brightlavender}{rgb}{0.75,0.58,0.89}
\definecolor{nudelavender}{rgb}{0.78, 0.66, 0.70}
\definecolor{codebg}{rgb}{0.10,0.10,0.14}
\definecolor{codekeyword}{rgb}{0.0,0.75,1.0}
\definecolor{codecomment}{rgb}{0.55,0.78,0.55}
\definecolor{codestring}{rgb}{0.95,0.65,0.40}

% --------------------------------------------------
% Dark Theme
% --------------------------------------------------
\pagecolor{black}
\color{white}

% Make table rules visible on a black background
\arrayrulecolor{brightlavender}

% --------------------------------------------------
% Section Style
% --------------------------------------------------
\titleformat{\section}
{\bfseries\Large\color{brightlavender}}
{\thesection}
{0.5em}
{}

\titleformat{\subsection}
{\bfseries\color{Thistle}}
{\thesubsection}
{0.5em}
{}

\titleformat{\subsubsection}
{\bfseries\color{nudelavender}}
{\thesubsubsection}
{0.5em}
{}

% --------------------------------------------------
% Bullet Style
% --------------------------------------------------
\setlist[itemize,1]{label=$\bullet$}
\setlist[itemize,2]{label=$\circ$}
\setlist[itemize,3]{label=$\bullet$}
\setlist[itemize,4]{label=$\circ$}

% --------------------------------------------------
% Code Listing Style (dark)
% --------------------------------------------------
\lstdefinestyle{darkcode}{
    backgroundcolor=\color{codebg},
    basicstyle=\ttfamily\small\color{white},
    keywordstyle=\color{codekeyword}\bfseries,
    commentstyle=\color{codecomment}\itshape,
    stringstyle=\color{codestring},
    numberstyle=\tiny\color{Gray},
    numbers=left,
    numbersep=8pt,
    breaklines=true,
    breakatwhitespace=false,
    frame=single,
    rulecolor=\color{brightlavender},
    showstringspaces=false,
    tabsize=4,
    columns=fullflexible,
    captionpos=b,
    aboveskip=10pt,
    belowskip=8pt,
}
\lstset{style=darkcode}

% --------------------------------------------------
% Title Style
% --------------------------------------------------
\title{\color{carminered}\textbf{Computer Networks}}
\author{\color{Thistle}Kiarash Sanei}
\date{\color{SkyBlue}June 2026}

\begin{document}

\maketitle

\begin{Large}
\tableofcontents
\end{Large}

\newpage

% ==================================================================
\section{Computer Networks and the Internet}
% ==================================================================

\subsection{The Internet: Two Views}

\subsubsection{The ``nuts-and-bolts'' view}
The Internet is a vast interconnection of computing devices and the
infrastructure that links them. It has three structural pieces:

\begin{itemize}
    \item \textbf{Hosts (end systems):} billions of connected devices
    running network applications at the Internet's ``edge'' --- laptops,
    phones, servers, but also cars, sensors, pacemakers, and ``smart''
    appliances (IoT).
    \item \textbf{Packet switches:} devices that forward \emph{packets}
    (chunks of data). The two main kinds are \textbf{routers} (network
    core) and \textbf{link-layer switches} (access/LAN).
    \item \textbf{Communication links:} fiber, copper, coaxial cable,
    radio, and satellite. Each link has a \textbf{transmission rate}
    (bandwidth) measured in bits/sec.
\end{itemize}

A \textbf{network} is a collection of devices, routers, and links managed
by a single organization. The Internet is therefore best described as a
\textbf{``network of networks''} --- a global mesh of interconnected
ISPs (Internet Service Providers).

\subsubsection{The ``services'' view}
Alternatively, the Internet is \emph{infrastructure that provides services
to distributed applications}: the Web, video streaming, email, games,
e-commerce, social media, and so on. It provides a \textbf{programming
interface} (the socket API) that lets applications ``hook into'' the
Internet's transport service --- much like a postal service offers
different delivery options.

\subsection{What Is a Protocol?}

A \textbf{protocol} defines the \emph{format} and \emph{order} of messages
exchanged between communicating entities, plus the \emph{actions} taken on
message transmission and receipt. Humans use protocols too (``Hi'' $\to$
``Hi'' $\to$ ``Got the time?'' $\to$ ``2:00''); on the Internet,
\emph{all} communication activity is governed by protocols, just between
machines instead of people.

\begin{itemize}
    \item Examples span every layer: HTTP (Web), TCP, IP, WiFi, 4/5G,
    Ethernet, and many more.
    \item Internet standards are published as \textbf{RFCs} (Request for
    Comments) by the \textbf{IETF} (Internet Engineering Task Force).
\end{itemize}

\subsection{The Network Edge: Hosts and Access Networks}

\textbf{Hosts} are split into \textbf{clients} (request service) and
\textbf{servers} (provide service, often housed in data centers). The
\textbf{access network} physically connects an end system to its first
(``edge'') router.

\begin{itemize}
    \item \textbf{Cable (HFC --- hybrid fiber coax):} shares a coaxial
    distribution network among homes; uses \textbf{frequency-division
    multiplexing (FDM)} so data and TV travel in different frequency
    bands. Asymmetric: up to roughly $40$ Mbps--$1.2$ Gbps downstream,
    $30$--$100$ Mbps upstream. The network device at the head end is the
    CMTS (Cable Modem Termination System).
    \item \textbf{DSL (Digital Subscriber Line):} uses the existing
    telephone line to the central office (DSLAM). \emph{Dedicated} (not
    shared) line: roughly $24$--$52$ Mbps down, $3.5$--$16$ Mbps up.
    Voice and data ride different frequencies on the same wire.
    \item \textbf{Home networks:} a single box typically combines cable/DSL
    modem, router, NAT, firewall, wired Ethernet ($1$ Gbps), and a WiFi
    access point ($54$/$450$ Mbps).
    \item \textbf{Enterprise networks:} a mix of Ethernet (100 Mbps / 1 / 10
    Gbps) and WiFi connecting switches and routers.
    \item \textbf{Wireless access:} WLAN/WiFi (802.11, $\sim$10s of meters)
    vs. wide-area cellular 4G/5G (operator-provided, $\sim$10s of km).
    \item \textbf{Data-center networks:} high-bandwidth links (10s--100s
    Gbps) connecting thousands of servers.
\end{itemize}

\subsubsection{Physical media}
\textbf{Guided} media carry signals in a solid: twisted pair (Cat 5/6
copper), coaxial cable, fiber optic (glass, light pulses; 10s--100s Gbps,
very low error rate, immune to electromagnetic noise). \textbf{Unguided}
media let signals propagate freely: radio. Radio link types include WiFi,
cellular, Bluetooth, terrestrial microwave, and satellite (geostationary
satellites add $\sim$270 ms end-to-end delay).

\subsection{The Network Core: Packet vs.\ Circuit Switching}

The core is a mesh of interconnected routers. There are two fundamental
ways to share core links among many application-level conversations.

\subsubsection{Circuit switching}
Resources (a ``circuit'') are \textbf{reserved} end-to-end for the duration
of a call.
\begin{enumerate}
    \item \textbf{Setup:} establish a circuit (e.g.\ when you dial a phone
    number); state is created in each switch along the path.
    \item \textbf{Communicate:} send data over the dedicated, guaranteed
    circuit.
    \item \textbf{Teardown:} release the circuit and remove the state.
\end{enumerate}
Resources are dedicated (no sharing $\Rightarrow$ guaranteed performance),
but a circuit sits \emph{idle} if the call has nothing to send. Two
multiplexing techniques:
\begin{itemize}
    \item \textbf{FDM:} the frequency spectrum is divided into bands; each
    call gets its own band.
    \item \textbf{TDM:} time is divided into frames, each frame into slots;
    each call gets a periodic slot and transmits at the full link rate
    \emph{during} its slot. Slot position within the frame identifies the
    conversation.
\end{itemize}

\subsubsection{Packet switching}
Used in the Internet. Messages are broken into \textbf{packets}, each
carrying a header (control info such as source/destination addresses).
\begin{itemize}
    \item \textbf{Store-and-forward:} a router must receive the \emph{entire}
    packet before it can begin transmitting it onto the next link.
    \item \textbf{No reserved capacity:} each packet is forwarded
    independently using the router's \textbf{forwarding table} (keyed on the
    destination address in the header).
    \item \textbf{Queues:} each output link has a buffer/queue; packets wait
    their turn. If a queue is full, arriving packets are \textbf{dropped}
    (lost).
\end{itemize}

\subsection{Two Key Network-Core Functions}

\begin{itemize}
    \item \textbf{Forwarding (``switching''):} a \emph{local} action ---
    move an arriving packet from an input link to the correct output link,
    consulting the local forwarding table.
    \item \textbf{Routing:} a \emph{global} action --- routing algorithms
    determine the end-to-end source--destination paths that packets take.
\end{itemize}

\subsection{Sharing the Network: Statistical Multiplexing}

Many \emph{flows} (streams of packets between two hosts) must share routers
and links. Two ways to allocate resources:
\begin{itemize}
    \item \textbf{Static allocation:} give each user a fixed share. Must
    provision for everyone's \emph{peak} demand $\Rightarrow$ wasteful.
    \item \textbf{Statistical multiplexing:} allocate dynamically, on
    demand. More efficient because
    \[
        \text{peak of aggregate demand} \;<\; \text{aggregate of peak demands.}
    \]
\end{itemize}

\noindent\textbf{Example.} Alice needs 10 units in the morning and 2 the
rest of the time; Bob needs 10 at night and 2 otherwise. Static: each gets
10 always $\Rightarrow$ provision 20. Dynamic: hand out 10 to whoever peaks
and 2 to the other $\Rightarrow$ provision only 12.

\medskip
\noindent The two canonical ways to implement multiplexing map directly onto
the two switching styles:
\begin{itemize}
    \item \textbf{Reservations via circuit switching:} explicitly reserve,
    use, then release resources (the ``restaurant reservation'' analogy).
    \item \textbf{Best-effort via packet switching:} just send packets and
    ``hope for the best'' (``first-come, first-served'').
\end{itemize}

\subsubsection{Why packet switching usually wins}
Consider a $1$ Gbps link with users that each need $100$ Mbps when active,
and are active only $10\%$ of the time.
\begin{itemize}
    \item \textbf{Circuit switching:} can support only $1000/100 = 10$ users
    (each statically reserved).
    \item \textbf{Packet switching:} with $35$ users, the probability that
    more than $10$ are simultaneously active is $< 0.0004$, so the link
    almost never overflows.
\end{itemize}
The per-user active count follows a \textbf{binomial distribution}: with $n$
independent users each active with probability $p$,
\[
    P(X = k) = \binom{n}{k}\, p^{k}\,(1-p)^{\,n-k}, \qquad
    P(X > 10) = \sum_{k=11}^{n} \binom{n}{k} p^{k}(1-p)^{\,n-k}.
\]
Packet switching is \textbf{great for bursty data} (resource sharing,
no call setup) but allows \textbf{congestion} (delay and loss from buffer
overflow), so it needs protocols for reliable transfer and congestion
control.

\subsection{Performance: Delay, Loss, Throughput}

\subsubsection{Four sources of nodal delay}
\[
    d_{\text{nodal}} = d_{\text{proc}} + d_{\text{queue}}
                     + d_{\text{trans}} + d_{\text{prop}}
\]
\begin{itemize}
    \item $d_{\text{proc}}$ --- \textbf{processing}: check bit errors,
    determine output link. Typically $<$ microseconds.
    \item $d_{\text{queue}}$ --- \textbf{queueing}: time waiting at the
    output link; depends on congestion.
    \item $d_{\text{trans}}$ --- \textbf{transmission}: time to push all $L$
    bits onto the link at rate $R$:
    \(\; d_{\text{trans}} = L/R.\)
    \item $d_{\text{prop}}$ --- \textbf{propagation}: time for a bit to
    travel the physical link of length $d$ at speed
    $s \approx 2\times 10^{8}$ m/s:
    \(\; d_{\text{prop}} = d/s.\)
\end{itemize}

\noindent\textbf{Transmission vs.\ propagation are different things.}
The caravan analogy: cars $=$ bits, a caravan $=$ a packet, a toll booth $=$
a link/router. ``Pushing'' the whole caravan through a booth is transmission
delay; a car driving between booths is propagation delay.

\subsubsection{Traffic intensity and queueing}
With average packet arrival rate $a$, packet length $L$, and link rate $R$,
the \textbf{traffic intensity} is $La/R$:
\[
    \frac{La}{R}\!\approx 0 \;\Rightarrow\; \text{small avg.\ queueing delay},
    \quad
    \frac{La}{R}\!\to 1 \;\Rightarrow\; \text{large delay},
    \quad
    \frac{La}{R} > 1 \;\Rightarrow\; \text{delay} \to \infty.
\]
\textbf{Rule of thumb:} design so that traffic intensity stays comfortably
below 1.

\subsubsection{Packet loss}
Router buffers are finite. When a packet arrives to a \emph{full} buffer it
is dropped. A lost packet may be retransmitted by the previous node, by the
source end system, or not at all.

\subsubsection{Throughput}
\textbf{Throughput} is the rate (bits/time) at which bits actually reach the
receiver (instantaneous vs.\ average). On a path with link rates
$R_s$ (server side) and $R_c$ (client side), end-to-end throughput is set by
the \textbf{bottleneck link}:
\[
    \text{throughput} = \min(R_s, R_c).
\]
If $N$ connections fairly share a backbone link of rate $R$, the per-connection
throughput is
\[
    \min\!\left(R_c,\; R_s,\; \frac{R}{N}\right),
\]
and in practice $R_c$ or $R_s$ (the access links) is usually the bottleneck.

\subsubsection{Traceroute}
\texttt{traceroute} measures delay to each router on the path: it sends three
probes with increasing TTL ($=i$) so probe $i$ expires at router $i$, which
returns an error; the sender times the round trip. \texttt{*} means no reply.

\subsection{Protocol Layers and the Internet Stack}

Networks are complex (hosts, routers, many media, applications). Layering
imposes structure and \textbf{modularity}: each layer relies on services from
the layer below and provides services to the layer above. A change inside one
layer (e.g.\ swapping a Layer-2 protocol) is transparent to the others ---
just as changing the airline gate procedure doesn't affect baggage handling.

\subsubsection{The five-layer Internet stack}
\begin{itemize}
    \item \textbf{Application} --- network applications and their protocols:
    HTTP, IMAP, SMTP, DNS. Unit: \emph{message}.
    \item \textbf{Transport} --- process-to-process data transfer:
    TCP, UDP. Unit: \emph{segment}.
    \item \textbf{Network} --- routing of datagrams source-to-destination:
    IP and routing protocols. Unit: \emph{datagram}.
    \item \textbf{Link} --- data transfer between neighboring nodes:
    Ethernet, 802.11 (WiFi), PPP. Unit: \emph{frame}.
    \item \textbf{Physical} --- bits ``on the wire.''
\end{itemize}
(The OSI model adds \emph{presentation} and \emph{session} layers, 6 and 5,
now effectively obsolete on the Internet.)

\subsubsection{The narrow waist}
\textbf{IP is the only protocol at the network layer} --- every host and
router speaks IP. This single ``narrow waist'' unifies the Internet and lets
arbitrary link-layer technologies below and arbitrary applications above
interoperate (the ``hourglass'').

\subsubsection{Encapsulation}
As a message descends the stack, each layer wraps the unit from above in its
own header (Matryoshka-doll style):
\[
    \underbrace{M}_{\text{message}}
    \;\to\; \underbrace{[H_t\,|\,M]}_{\text{segment}}
    \;\to\; \underbrace{[H_n\,|\,H_t\,|\,M]}_{\text{datagram}}
    \;\to\; \underbrace{[H_l\,|\,H_n\,|\,H_t\,|\,M]}_{\text{frame}}.
\]
At the destination, headers are \emph{peeled} off going up. Key insight from
the ``postal/CEO secretary'' analogy: \textbf{each layer only reads the header
written by its peer at the same layer.} Layer-3 (IP) destination is always the
final endpoint; the Layer-2 header changes hop-by-hop.

\subsubsection{Who implements which layers}
\begin{itemize}
    \item \textbf{End hosts} implement \emph{all} layers (1--7). In hardware
    (the NIC): layers 1--2. In the OS: layers 3--4. In applications: layer 7.
    \item \textbf{Routers} implement only layers 1--3: they parse and forward,
    rewriting the Layer-2 header for the next hop, but do not provide reliable
    delivery (4) or touch application data (7).
\end{itemize}

\subsubsection{Layer 3 is best-effort}
The network layer offers a \textbf{best-effort} service model: packets are
size-limited and may be lost, reordered, or corrupted; the network ``tries its
best'' but gives no guarantee and won't even tell you when delivery fails.
Higher layers (transport) must add reliability if the application needs it.

\subsection{Network Security (Introduction)}

The Internet was originally designed for ``a group of mutually trusting users
attached to a transparent network,'' so security has always been retrofitted.

\begin{itemize}
    \item \textbf{Packet sniffing:} a promiscuous NIC on a broadcast medium
    (shared Ethernet, WiFi) records all passing packets --- including
    cleartext passwords. (Wireshark is a benign sniffer used in labs.)
    \item \textbf{IP spoofing:} inject packets with a forged source address.
    \item \textbf{Denial of Service (DoS):} overwhelm a target (server,
    bandwidth) with bogus traffic, often from a \emph{botnet} of compromised
    hosts (DDoS).
\end{itemize}

\noindent\textbf{Lines of defense:} authentication (e.g.\ SIM hardware
identity), confidentiality (encryption), integrity (digital signatures),
access restrictions (VPNs, passwords), and \textbf{firewalls}
(``off-by-default'' middleboxes that filter traffic and react to DoS).

\subsection{A Brief History}

\begin{itemize}
    \item \textbf{1950s--60s, Cold War:} the US wanted a communications
    system that could survive a nuclear attack. \textbf{ARPA} (1958, a
    response to Sputnik; later DARPA) funded large-scale networking research.
    \item \textbf{Paul Baran (RAND, early 1960s):} survivable comms via
    \emph{decentralized architecture} (no single point of failure),
    \emph{message blocks} (packets sent independently and reassembled), and
    \emph{hot-potato routing} (forward as fast as possible).
    \item \textbf{Donald Davies (UK):} coined the term ``packet.''
    \textbf{Leonard Kleinrock (MIT, 1961--62):} packet-switching theory.
    \textbf{Licklider:} vision of a ``galactic network.''
    \item \textbf{Kleinrock $\to$ Roberts:} Roberts wrote the ARPANET plan
    (1967); DARPA's 1968 RFQ for the packet switches (Interface Message
    Processors, IMPs) was won by \textbf{BBN}.
    \item \textbf{1969:} first nodes --- UCLA (Kleinrock's Network Measurement
    Center) and SRI. The first host-to-host message (``LO\dots'' then the
    system crashed) was sent UCLA$\to$SRI; UCSB and Utah followed.
    \item \textbf{1970:} the Network Working Group (Crocker) finished NCP,
    the first host-to-host protocol. \textbf{1972:} Kahn demoed ARPANET at
    ICCC; Ray Tomlinson (BBN) wrote email.
    \item \textbf{Open-architecture / internetworking:} Kahn's four ground
    rules --- (1) each network stands on its own, (2) best-effort delivery
    with source retransmission, (3) simple stateless ``black boxes''
    (gateways/routers), (4) no global control. Kahn \& Cerf's 1974 paper
    defined \textbf{TCP/IP}, later split into IP (addressing/forwarding),
    TCP (reliability), and UDP (direct access to IP).
    \item \textbf{Documentation:} Crocker started the \textbf{RFC} series
    (1969); the \textbf{IETF} (founded 1986) makes voluntary standards
    published as RFCs.
\end{itemize}

% ==================================================================
\section{The Application Layer}
% ==================================================================

\subsection{Principles of Network Applications}

You write applications that run on \emph{end systems} and communicate over
the network; you never need to write software for core devices (they don't
run user applications). Two dominant architectures:

\begin{itemize}
    \item \textbf{Client--server:} an always-on server with a permanent IP
    address (often in a data center for scaling) serves intermittently
    connected clients. Clients don't talk to each other directly.
    Examples: HTTP, IMAP, FTP.
    \item \textbf{Peer-to-peer (P2P):} no always-on server; peers request and
    provide service to each other. \emph{Self-scaling} (each new peer adds
    capacity \emph{and} demand) but harder to manage (peers churn and change
    IPs). Example: BitTorrent.
\end{itemize}

\subsubsection{Processes, sockets, and addressing}
A \textbf{process} is a running program. Processes on different hosts
communicate by exchanging messages through \textbf{sockets} --- the
``door'' between the application and the transport layer. To be addressable,
a process needs an \textbf{IP address} (identifies the host) \emph{and} a
\textbf{port number} (identifies the process/socket on that host), because
many processes run on one host. Example well-known ports: HTTP $=80$,
SMTP $=25$, DNS $=53$.

\subsubsection{What an application-layer protocol defines}
The types of messages (request/response), message \emph{syntax} (fields and
how they're delineated), message \emph{semantics} (meaning of fields), and
\emph{rules} for when/how to send and respond. \textbf{Open} protocols are
defined in RFCs (HTTP, SMTP) and enable interoperability;
\textbf{proprietary} protocols are closed (e.g.\ some video-conferencing
codecs).

\subsection{What Transport Service Does an App Need?}

Apps care about four dimensions:
\begin{itemize}
    \item \textbf{Data integrity:} file transfer needs 100\% reliability;
    audio tolerates some loss.
    \item \textbf{Timing:} interactive games / telephony need low delay.
    \item \textbf{Throughput:} multimedia needs a minimum rate;
    ``elastic'' apps use whatever they get.
    \item \textbf{Security:} encryption, integrity, authentication.
\end{itemize}

\subsubsection{TCP vs.\ UDP at a glance}
\begin{itemize}
    \item \textbf{TCP:} reliable, in-order delivery; flow control;
    congestion control; connection setup. Does \emph{not} guarantee timing
    or minimum throughput.
    \item \textbf{UDP:} unreliable, unordered; no setup, no flow/congestion
    control. Minimal overhead. ``Why bother?'' --- it's fast and works under
    congestion, and apps can add their own reliability on top (HTTP/3 does
    exactly this).
\end{itemize}

\subsubsection{Securing TCP: TLS}
Vanilla TCP/UDP sockets send cleartext. \textbf{TLS} (Transport Layer
Security) is implemented in the \emph{application layer} (apps link a TLS
library that uses TCP underneath); it provides encrypted connections, data
integrity, and endpoint authentication.

\subsection{Socket Programming (Completing the Roadmap)}

The slides list socket programming with UDP and TCP but don't show code.
Here are minimal Python examples that make the client--server paradigm
concrete.

\subsubsection{UDP: connectionless}
There is no handshake; every datagram is self-contained and must carry the
destination address.

\begin{lstlisting}[language=Python, caption={UDP echo server and client}]
# --- UDP server ---
from socket import socket, AF_INET, SOCK_DGRAM

server_port = 12000
s = socket(AF_INET, SOCK_DGRAM)        # SOCK_DGRAM = UDP
s.bind(('', server_port))              # listen on all interfaces
print("UDP server ready")

while True:
    message, client_addr = s.recvfrom(2048)   # blocks for a datagram
    reply = message.decode().upper()
    s.sendto(reply.encode(), client_addr)     # must specify dest each time

# --- UDP client ---
from socket import socket, AF_INET, SOCK_DGRAM

server_name, server_port = 'localhost', 12000
s = socket(AF_INET, SOCK_DGRAM)
msg = input("Input lowercase sentence: ")
s.sendto(msg.encode(), (server_name, server_port))
reply, _ = s.recvfrom(2048)
print(reply.decode())
s.close()
\end{lstlisting}

\subsubsection{TCP: connection-oriented}
The server has a \emph{welcoming} socket; \texttt{accept()} returns a
\emph{new} connection socket dedicated to that one client (this is the
4-tuple demultiplexing we'll see in the transport layer).

\begin{lstlisting}[language=Python, caption={TCP echo server and client}]
# --- TCP server ---
from socket import socket, AF_INET, SOCK_STREAM

server_port = 12000
welcome = socket(AF_INET, SOCK_STREAM)    # SOCK_STREAM = TCP
welcome.bind(('', server_port))
welcome.listen(1)                         # backlog of pending connections
print("TCP server ready")

while True:
    conn, addr = welcome.accept()         # new socket per client
    sentence = conn.recv(1024).decode()
    conn.send(sentence.upper().encode())
    conn.close()

# --- TCP client ---
from socket import socket, AF_INET, SOCK_STREAM

server_name, server_port = 'localhost', 12000
s = socket(AF_INET, SOCK_STREAM)
s.connect((server_name, server_port))     # triggers the 3-way handshake
s.send(input("Input sentence: ").encode())
print(s.recv(1024).decode())
s.close()
\end{lstlisting}

\subsection{Web and HTTP}

A \textbf{web page} consists of objects (an HTML base file plus referenced
objects --- JPEGs, scripts, etc.), each addressable by a \textbf{URL} of the
form \texttt{host\-name/path\-name}. \textbf{HTTP} (HyperText Transfer
Protocol) is the Web's application-layer protocol, developed by Tim
Berners-Lee at CERN in 1989 (HTTP/0.9 in 1991, HTTP/1.0 in 1996, HTTP/1.1 in
1997).

\begin{itemize}
    \item \textbf{Client--server, over TCP:} client opens a TCP connection to
    port 80, exchanges request/response messages, then the connection closes.
    \item \textbf{Stateless:} the server keeps \emph{no} memory of past client
    requests. Stateful protocols are complex --- after a crash, the two sides'
    views of ``state'' can diverge and must be reconciled.
\end{itemize}

\subsubsection{Non-persistent vs.\ persistent HTTP}
\begin{itemize}
    \item \textbf{Non-persistent (HTTP/1.0):} at most one object per TCP
    connection. Each object costs roughly $2\,RTT$ (one RTT to set up TCP,
    one for the request + first bytes) plus the file's transmission time.
    \[
        T_{\text{non-persistent (per object)}} = 2\,RTT + T_{\text{transmit}}.
    \]
    \item \textbf{Persistent (HTTP/1.1):} the server leaves the connection
    open; multiple objects reuse it. As little as one RTT for all referenced
    objects.
\end{itemize}

\subsubsection{HTTP message format}
Requests are ASCII (human-readable). A request line (method, URL, version)
is followed by header lines; methods include \texttt{GET}, \texttt{POST}
(form data in the entity body), \texttt{HEAD} (headers only), and
\texttt{PUT} (upload/replace).

\begin{lstlisting}[caption={A raw HTTP/1.1 GET request (CRLF = \textbackslash r\textbackslash n)}]
GET /index.html HTTP/1.1
Host: www-net.cs.umass.edu
User-Agent: Mozilla/5.0 (...) Firefox/80.0
Accept: text/html,application/xhtml+xml
Accept-Language: en-us,en;q=0.5
Accept-Encoding: gzip,deflate
Connection: keep-alive
                         <-- a blank line (just CRLF) ends the headers
\end{lstlisting}

A response begins with a \textbf{status line} (version, status code,
phrase), then headers, then the data. Common status codes:
\textbf{200} OK, \textbf{301} Moved Permanently (new URL in
\texttt{Location:}), \textbf{400} Bad Request, \textbf{404} Not Found,
\textbf{505} HTTP Version Not Supported.

\subsubsection{Trying HTTP by hand}
You can speak HTTP manually with \texttt{netcat}/\texttt{ncat}:

\begin{lstlisting}[language=bash, caption={Manual HTTP request over a raw TCP connection}]
# Open a TCP connection to port 80, then type the request:
nc -c -v gaia.cs.umass.edu 80        # macOS
ncat -C gaia.cs.umass.edu 80         # Windows

GET /kurose_ross/interactive/index.php HTTP/1.1
Host: gaia.cs.umass.edu
                                     # press Enter twice (blank line)
\end{lstlisting}

\noindent A more typical programmatic fetch:

\begin{lstlisting}[language=Python, caption={A raw HTTP GET via a socket}]
from socket import socket, AF_INET, SOCK_STREAM

host = 'gaia.cs.umass.edu'
s = socket(AF_INET, SOCK_STREAM)
s.connect((host, 80))
request = (f"GET / HTTP/1.1\r\n"
           f"Host: {host}\r\n"
           f"Connection: close\r\n\r\n")   # blank line terminates headers
s.send(request.encode())

response = b""
while True:
    chunk = s.recv(4096)
    if not chunk:
        break
    response += chunk
s.close()
print(response.decode(errors='replace')[:500])
\end{lstlisting}

\subsection{Maintaining State: Cookies}

HTTP itself is stateless, but sites keep state with \textbf{cookies}, which
have four components:
\begin{enumerate}
    \item A \texttt{Set-Cookie:} header in the HTTP \emph{response}.
    \item A \texttt{Cookie:} header in subsequent HTTP \emph{requests}.
    \item A cookie file on the user's host, managed by the browser.
    \item A back-end database at the site.
\end{enumerate}
On a first visit the site creates a unique ID and a database entry; every
later request carries the ID so the site can ``recognize'' the user.

\begin{itemize}
    \item \textbf{Uses:} authorization, shopping carts, recommendations,
    web-email session state.
    \item \textbf{First-party} cookie: from the site you chose to visit.
    \textbf{Third-party / tracking} cookie: from an ad/tracker
    (e.g.\ \texttt{AdX.com}) embedded across many sites, letting a common ID
    follow you around --- often via an invisible link rather than a visible
    ad. Increasingly disabled by default and regulated (e.g.\ EU GDPR treats
    identifying cookies as personal data).
\end{itemize}

\subsection{Web Caching (Proxy Servers)}

A \textbf{web cache} sits between clients and origin servers, acting as both
server (to the client) and client (to the origin). The browser sends all
requests to the cache; on a hit it returns the object directly, on a miss it
fetches from the origin, caches a copy, and returns it.

\textbf{Why cache?} Reduce client response time (cache is nearer), reduce
traffic on the institution's access link, and let ``poor'' content providers
deliver effectively.

\subsubsection{Worked example}
Origin RTT $=2$ s, object size $=100$ Kbits, request rate $=15$/s, access
link $=1.54$ Mbps, LAN $=1$ Gbps. Average data rate to browsers
$=15 \times 100\text{ Kbits} = 1.50$ Mbps.
\begin{itemize}
    \item \textbf{No cache:} access-link utilization
    $=1.50/1.54 \approx 0.97 \Rightarrow$ huge queueing delay
    (end-to-end delay $=$ 2 s $+$ \emph{minutes} $+ \mu$s).
    \item \textbf{Faster link (154 Mbps):} utilization $\approx 0.0097$,
    delay in msec --- but expensive.
    \item \textbf{Install a cache (hit rate $0.4$):} only $60\%$ of traffic
    hits the access link $\Rightarrow 0.6 \times 1.50 = 0.9$ Mbps,
    utilization $0.9/1.54 \approx 0.58 \Rightarrow$ low delay. Average
    end-to-end delay $\approx 0.6(2.01) + 0.4(\sim\!\text{msec})
    \approx 1.2$ s --- lower \emph{and} cheaper than the fast link.
\end{itemize}

\subsubsection{Conditional GET}
To avoid sending an object the browser already has fresh, the client adds
\texttt{If-Modified-Since: <date>}. The server replies
\texttt{304 Not Modified} (no body) if unchanged, or \texttt{200 OK} with the
object if it changed.

\subsection{HTTP/2 and HTTP/3}

\textbf{HTTP/2} (RFC 7540, 2015) targets reduced delay for multi-object
pages. HTTP/1.1 introduced pipelined GETs over one TCP connection, but the
server replies FCFS, so a small object can suffer \textbf{head-of-line (HOL)
blocking} behind a large one, and TCP loss recovery stalls everything.
HTTP/2 keeps methods/status codes/most headers but adds:
\begin{itemize}
    \item client-specified object \textbf{priority} (not necessarily FCFS),
    \item server \textbf{push} of unrequested objects,
    \item dividing objects into \textbf{frames} and \emph{interleaving} them
    to mitigate HOL blocking.
\end{itemize}

\textbf{HTTP/3} addresses HTTP/2's remaining problem: a single TCP connection
means one packet loss still stalls all streams. HTTP/3 runs over
\textbf{UDP} (via QUIC), adding security plus \emph{per-object} error and
congestion control.

\subsection{Email: SMTP and IMAP (Completing the Roadmap)}

The chapter roadmap lists ``E-mail, SMTP, IMAP'' but the slides skip it.
Here is the essential picture.

\begin{itemize}
    \item \textbf{Components:} user agents (mail clients), mail servers
    (each with an outgoing \emph{message queue} and per-user \emph{mailbox}),
    and the \textbf{SMTP} protocol between servers.
    \item \textbf{SMTP (Simple Mail Transfer Protocol, RFC 5321):} a
    \emph{push} protocol over TCP \textbf{port 25}; uses persistent
    connections; commands and replies are 7-bit ASCII; a message ends with a
    line containing only ``\texttt{.}''.
    \item \textbf{Retrieval:} SMTP delivers mail \emph{to} the recipient's
    server, but a user agent \emph{pulls} mail from its mailbox with
    \textbf{IMAP} (keeps messages on the server, supports folders/search) or
    the older POP3. Webmail uses HTTP for this last hop.
\end{itemize}

\begin{lstlisting}[language=bash, caption={A bare SMTP dialog (client lines after the prompts)}]
telnet mail.example.com 25
S: 220 mail.example.com ESMTP ready
C: HELO relay.local
S: 250 Hello relay.local
C: MAIL FROM:<alice@example.com>
S: 250 OK
C: RCPT TO:<bob@example.com>
S: 250 OK
C: DATA
S: 354 End data with <CRLF>.<CRLF>
C: Subject: Hi Bob
C:
C: Your days are numbered.
C: .
S: 250 OK queued
C: QUIT
S: 221 Bye
\end{lstlisting}

\subsection{The Domain Name System (DNS)}

\textbf{DNS} translates human-readable domain names (\texttt{www.google.com})
into IP addresses (\texttt{74.125.25.99}). It is itself an
\emph{application-layer} protocol, even though name resolution is a core
Internet function --- complexity is pushed to the network's edge.

\subsubsection{Goals and why it's not centralized}
DNS must be \textbf{scalable} (billions of records, trillions of
queries/day), \textbf{highly available} (no single point of failure),
and \textbf{fast} (almost every connection begins with a lookup, so a slow
DNS makes everything slow). Centralizing would fail on all of these
(single point of failure, traffic volume, distance, maintenance). It is
read-heavy: far more lookups than updates.

\subsubsection{The hierarchy}
DNS arranges \textbf{name servers} in a tree; if one can't answer, it
\emph{delegates} to a child.
\begin{itemize}
    \item \textbf{Root servers} (responsible for everything; only 13
    root-server \emph{names}, \texttt{a}--\texttt{m.root-servers.net}, but
    thousands of physical servers via \textbf{anycast}).
    \item \textbf{Top-Level Domain (TLD) servers:} \texttt{.com},
    \texttt{.org}, \texttt{.edu}, country domains, etc.\ (Verisign/Network
    Solutions for \texttt{.com}; Educause for \texttt{.edu}). Over 1500 TLDs
    today.
    \item \textbf{Authoritative servers:} an organization's own servers with
    the definitive records for its hosts.
\end{itemize}

\textbf{Three kinds of hierarchy:} names (dotted), authority (a
\emph{zone} = an administrative authority that can delegate sub-zones), and
infrastructure (each server knows only a subset).

\subsubsection{Iterative vs.\ recursive resolution}
\begin{itemize}
    \item \textbf{Iterative:} the contacted server replies ``I don't know,
    but ask \emph{that} server.'' The client walks the tree itself.
    \item \textbf{Recursive:} the client offloads the whole job to its
    \textbf{recursive resolver} (e.g.\ ISP-run, or \texttt{8.8.8.8}
    /Google, \texttt{1.1.1.1}/Cloudflare). The resolver does the iterative
    queries and returns the final answer. Resolvers build large shared caches
    and reduce load on name servers. Your machine's \textbf{stub resolver}
    just asks the recursive resolver.
\end{itemize}
Answers are \textbf{cached} for a \textbf{TTL} (e.g.\ 86\,400 s $=$ 24 h).

\subsubsection{Packet format and records}
DNS runs over \textbf{UDP port 53} (lightweight, no handshake, no
per-connection state; usually fits one packet). Large transfers (e.g.\ zone
transfers to secondary servers) fall back to TCP. The header carries a 16-bit
\textbf{ID} (matches responses to queries; should be random for security),
\textbf{flags} (QR query/response bit; RD recursion-desired bit), and four
record \textbf{counts}. The payload holds \textbf{resource records (RRs)}
--- name--value pairs sorted into Question / Answer / Authority / Additional
sections.

Record types you should know:
\begin{itemize}
    \item \textbf{A} --- name $\to$ IPv4 address.
    \item \textbf{AAAA} --- name $\to$ IPv6 address.
    \item \textbf{NS} --- delegates a domain to another name server.
    \item \textbf{CNAME} --- alias for a canonical name.
    \item \textbf{MX} --- mail server for a domain (with a priority; try
    lowest number first).
\end{itemize}

\subsubsection{Doing it yourself with \texttt{dig}}
\begin{lstlisting}[language=bash, caption={Walking the hierarchy manually with +norecurse}]
# Ask the root directly; it delegates to .edu
$ dig +norecurse eecs.berkeley.edu @198.41.0.4
;; AUTHORITY SECTION:
edu.  172800  IN  NS  a.edu-servers.net.
;; ADDITIONAL SECTION:
a.edu-servers.net.  172800  IN  A  192.5.6.30   # glue record

# Now ask .edu; it delegates to berkeley.edu
$ dig +norecurse eecs.berkeley.edu @192.5.6.30

# Finally ask berkeley.edu; it answers authoritatively
$ dig +norecurse eecs.berkeley.edu @128.32.136.3
;; ANSWER SECTION:
eecs.berkeley.edu.  86400  IN  A  23.185.0.1
\end{lstlisting}

\begin{lstlisting}[language=Python, caption={Resolving names from a program}]
import socket
print(socket.gethostbyname('www.google.com'))      # -> e.g. 142.250.x.x
# Modern, supports IPv6 and more:
for fam, *_ , sockaddr in socket.getaddrinfo('example.com', 443):
    print(sockaddr[0])
\end{lstlisting}

\subsubsection{Other uses of DNS}
\begin{itemize}
    \item \textbf{Email} via MX records.
    \item \textbf{Load balancing:} return multiple A records (client picks
    one; server shuffles order) for coarse-grained balancing and resiliency.
    \item \textbf{Geographic load balancing:} return different answers based
    on the (resolver's) location to steer clients to a ``nearby'' server ---
    inherently approximate.
\end{itemize}

\subsubsection{DNS security}
\textbf{DDoS} on root servers (mitigated by traffic filtering and resolver
caching of TLD addresses) or the more dangerous TLD servers;
\textbf{spoofing} / \textbf{cache poisoning} (returning bogus replies).
\textbf{DNSSEC} (RFC 4033) adds authentication.

\subsection{P2P and BitTorrent}

\subsubsection{File-distribution time: client--server vs.\ P2P}
Distributing a file of size $F$ to $N$ peers; server upload $u_s$, peer $i$
upload/download $u_i, d_i$, with $d_{\min}=\min_i d_i$.

\[
    D_{\text{c-s}} \;\geq\; \max\!\left\{\frac{N F}{u_s},\; \frac{F}{d_{\min}}\right\}
    \qquad\text{(grows \emph{linearly} in } N\text{)}
\]
\[
    D_{\text{P2P}} \;\geq\; \max\!\left\{\frac{F}{u_s},\; \frac{F}{d_{\min}},\;
        \frac{N F}{\,u_s + \sum_i u_i\,}\right\}
\]
In P2P, demand $N F$ \emph{and} supply $u_s + \sum_i u_i$ both grow with $N$,
so distribution time scales far better --- the source of P2P's
\textbf{self-scalability}.

\subsubsection{BitTorrent mechanics}
\begin{itemize}
    \item A file is split into \textbf{256-Kb chunks}; a \textbf{torrent} is
    the group of peers exchanging chunks; a \textbf{tracker} tracks
    participating peers.
    \item A new peer registers with the tracker, gets a list of peers,
    connects to a subset (``neighbors''), and accumulates chunks while
    uploading what it has. \textbf{Churn:} peers join/leave; a finished peer
    may leave (selfish) or seed (altruistic).
    \item \textbf{Requesting --- rarest first:} Alice periodically asks
    neighbors for their chunk lists and requests the \emph{rarest} missing
    chunks first.
    \item \textbf{Sending --- tit-for-tat:} Alice sends to the four peers
    feeding her fastest (others are ``choked''); re-evaluated every 10 s.
    Every 30 s she ``optimistically unchokes'' a random peer to discover
    better partners. Higher upload rate $\Rightarrow$ better trading
    partners $\Rightarrow$ faster downloads.
\end{itemize}

\subsection{Video Streaming and CDNs}

Streaming video (Netflix, YouTube, Amazon Prime) was $\sim$80\% of
residential ISP traffic (2020). Challenges: \textbf{scale} ($\sim$1B users)
and \textbf{heterogeneity} (wired vs.\ mobile, rich vs.\ poor bandwidth).

\subsubsection{Video basics}
A video is a sequence of images (e.g.\ 24--30 fps); each image is an array of
pixels. \textbf{Coding} exploits redundancy to cut bits: \emph{spatial}
(within an image --- e.g.\ run-length encode a block of identical color) and
\emph{temporal} (between frames --- send only the difference from the
previous frame). \textbf{CBR} (constant bit rate) vs.\ \textbf{VBR} (variable
bit rate). Examples: MPEG-1 $1.5$ Mbps, MPEG-2 $3$--$6$ Mbps, MPEG-4
$64$ Kbps--$12$ Mbps.

\subsubsection{Playout buffering}
Network delay is variable (\textbf{jitter}), but playout must be continuous.
A \textbf{client-side buffer} absorbs jitter: the client pre-fills the buffer
and plays out at a constant rate while reception varies, adding a
\textbf{playout delay}.

\subsubsection{DASH (Dynamic Adaptive Streaming over HTTP)}
\begin{itemize}
    \item \textbf{Server:} divides the video into \emph{chunks}, encodes each
    at multiple rates stored in separate files, and provides a
    \textbf{manifest} listing chunk URLs; files are replicated across CDN
    nodes.
    \item \textbf{Client (the ``intelligence''):} periodically estimates
    available bandwidth and, consulting the manifest, requests one chunk at a
    time --- choosing \emph{when} (avoid buffer starvation/overflow),
    \emph{what rate} (highest sustainable quality), and \emph{where} (a CDN
    node that is close or uncongested).
\end{itemize}

\subsubsection{Content Distribution Networks}
A single ``megaserver'' doesn't scale (single point of failure, congestion,
long paths). A \textbf{CDN} stores copies at many geographically distributed
sites. Two philosophies:
\begin{itemize}
    \item \textbf{Enter deep} (e.g.\ Akamai, $\sim$240k servers in $>$120
    countries): push servers \emph{deep} into access networks, close to
    users.
    \item \textbf{Bring home} (e.g.\ Limelight): fewer, larger clusters in
    points-of-presence (POPs) near access nets.
\end{itemize}
\textbf{Netflix OpenConnect} stores content at CDN nodes; the client gets a
manifest and pulls each chunk at the highest supportable rate (or a different
copy if a path is congested). Core \textbf{OTT} (``over the top'') challenges:
\emph{which} content goes in \emph{which} node, and \emph{from where / at what
rate} to retrieve it.

% ==================================================================
\section{The Transport Layer}
% ==================================================================

\subsection{Services and Principles}

The transport layer provides \textbf{logical end-to-end communication between
processes}, building on the network layer's host-to-host service.
\begin{itemize}
    \item \textbf{Sender:} breaks application messages into
    \textbf{segments}, passes them to the network layer.
    \item \textbf{Receiver:} reassembles segments into messages, passes them
    up to the application.
\end{itemize}

\noindent\textbf{Household analogy:} 12 kids in Ann's house writing to 12 kids
in Bill's house. Houses $=$ hosts, kids $=$ processes, letters $=$ messages,
the postal service $=$ network layer, and \emph{Ann and Bill}
(who sort mail to the right sibling) $=$ the transport protocol.

\subsection{Multiplexing and Demultiplexing}

\begin{itemize}
    \item \textbf{Multiplexing (sender):} gather data from multiple sockets,
    add transport headers.
    \item \textbf{Demultiplexing (receiver):} use header fields to deliver
    each segment to the correct socket.
\end{itemize}

\subsubsection{Connectionless (UDP) demux}
A UDP socket is identified by the \textbf{2-tuple}
(\emph{dest IP, dest port}). Two datagrams with the same destination port but
different source IPs/ports are delivered to the \emph{same} socket.

\subsubsection{Connection-oriented (TCP) demux}
A TCP socket is identified by the full \textbf{4-tuple}
(\emph{src IP, src port, dest IP, dest port}). A server therefore supports
many simultaneous sockets, one per connecting client --- which is exactly why
\texttt{accept()} returns a fresh socket. Three segments all sent to
(B, 80) but from different sources demux to different sockets.

\subsection{UDP: User Datagram Protocol (RFC 768)}

``No frills, bare bones, best effort.'' Segments may be lost or delivered
out of order. Benefits: \textbf{no connection setup} (no RTT penalty),
\emph{no} connection state, \textbf{small header}, \emph{no} congestion
control (can blast as fast as desired; functions under congestion).
\textbf{Uses:} DNS, SNMP, streaming multimedia (loss-tolerant,
rate-sensitive), and HTTP/3 (which re-adds reliability and congestion control
at the application layer).

\subsubsection{UDP segment and checksum}
Header (8 bytes): source port, dest port, length (bytes, incl.\ header),
checksum. The \textbf{Internet checksum} detects flipped bits:

\begin{itemize}
    \item \textbf{Sender:} treat the segment (plus a pseudo-header with IP
    addresses) as a sequence of 16-bit integers; compute the one's-complement
    sum (a carry out of the top bit is added back --- ``wraparound''); store
    the one's complement of that sum in the checksum field.
    \item \textbf{Receiver:} sum all 16-bit words \emph{including} the
    checksum; the result should be all ones. Not equal $\Rightarrow$ error
    detected. (It's \textbf{weak} protection: compensating bit flips can leave
    the checksum unchanged.)
\end{itemize}

\begin{lstlisting}[language=Python, caption={The 16-bit Internet checksum}]
def internet_checksum(data: bytes) -> int:
    if len(data) % 2:                      # pad to an even number of bytes
        data += b'\x00'
    total = 0
    for i in range(0, len(data), 2):
        word = (data[i] << 8) + data[i + 1]
        total += word
        total = (total & 0xFFFF) + (total >> 16)   # fold the carry (wraparound)
    return ~total & 0xFFFF                  # one's complement

# Verify: summing data + checksum should give 0xFFFF (all ones)
\end{lstlisting}

\subsection{Principles of Reliable Data Transfer (rdt)}

Goal: provide a \textbf{reliable channel} abstraction to the application on
top of an \emph{unreliable} network channel. Complexity depends on what the
channel can do to packets (lose / corrupt / reorder). Neither side knows the
other's state unless told via a message --- ``that's why we need a protocol.''
We develop the sender and receiver as \textbf{finite state machines (FSMs)},
considering unidirectional data (control info still flows both ways).

\begin{itemize}
    \item \textbf{rdt 1.0:} perfectly reliable channel --- just send and
    receive.
    \item \textbf{rdt 2.0:} channel may \emph{flip bits}. Add a
    \emph{checksum} to detect errors and \textbf{ACK}/\textbf{NAK} feedback:
    the sender retransmits on a NAK (``stop-and-wait''). \emph{Fatal flaw:}
    what if the ACK/NAK itself is corrupted?
    \item \textbf{rdt 2.1:} add a 1-bit \textbf{sequence number} (0/1) so the
    receiver can discard duplicates; the sender retransmits if the ACK/NAK is
    corrupted. Two sequence numbers suffice for stop-and-wait.
    \item \textbf{rdt 2.2:} a \textbf{NAK-free} protocol --- the receiver
    ACKs the last correctly received packet (including its seq \#); a
    \emph{duplicate ACK} means ``retransmit.'' (TCP uses this idea.)
    \item \textbf{rdt 3.0:} channel may also \emph{lose} packets. Add a
    \textbf{countdown timer}: if no ACK arrives in a ``reasonable'' time,
    retransmit. A delayed (not lost) packet causes a harmless duplicate that
    sequence numbers already handle.
\end{itemize}

\subsubsection{Why stop-and-wait performs terribly}
With a $1$ Gbps link, $15$ ms one-way prop.\ delay (so $RTT = 30$ ms), and
$L = 8000$-bit packets: $d_{\text{trans}} = L/R = 8\,\mu$s. Sender utilization:
\[
    U_{\text{sender}} = \frac{L/R}{RTT + L/R}
    = \frac{0.008}{30.008} \approx 0.00027.
\]
The protocol --- not the link --- limits performance.

\subsection{Pipelined Protocols}

\textbf{Pipelining:} allow multiple ``in-flight,'' yet-to-be-ACKed packets.
This needs a larger range of sequence numbers and buffering at sender and/or
receiver. Three packets in flight improves utilization 3$\times$:
\[
    U_{\text{sender}} = \frac{3\,L/R}{RTT + L/R}.
\]

\subsubsection{Go-Back-N (GBN)}
\begin{itemize}
    \item Sender keeps a window of up to $N$ consecutive transmitted-but-unACKed
    packets, with a $k$-bit sequence number.
    \item \textbf{Cumulative ACK:} ACK($n$) acknowledges \emph{all} packets up
    to and including $n$; on receipt the window slides to begin at $n+1$.
    \item A single timer covers the oldest in-flight packet;
    \textbf{timeout($n$)} retransmits packet $n$ and \emph{all} higher-seq
    packets in the window.
    \item \textbf{Receiver} is simple: it only tracks the next expected seq
    (\texttt{rcv\_base}), discards out-of-order packets, and re-ACKs the
    highest in-order seq received.
\end{itemize}

\subsubsection{Selective Repeat (SR)}
\begin{itemize}
    \item The receiver \emph{individually} ACKs each correctly received packet
    and \emph{buffers} out-of-order packets for later in-order delivery.
    \item The sender keeps a (conceptual) timer \emph{per} unACKed packet and
    on timeout retransmits \emph{only} that one packet.
    \item \textbf{The dilemma:} sequence numbers wrap around. To avoid the
    receiver mistaking a retransmission for new data, the
    \textbf{window size must be} $\leq \tfrac{1}{2}$ the sequence-number
    space. (Receiver windows in $[\texttt{rcvbase}, \texttt{rcvbase}+N-1]$ and
    $[\texttt{rcvbase}-N, \texttt{rcvbase}-1]$ must not overlap.)
\end{itemize}

\subsection{TCP: Connection-Oriented Transport (RFCs 793, 1122, 2018, 5681, 7323)}

TCP is point-to-point, reliable, in-order \textbf{byte stream}
(no message boundaries), \textbf{full-duplex}, connection-oriented
(handshake first), flow-controlled, and congestion-controlled. It pipelines
with \textbf{cumulative ACKs} and a window set by flow and congestion
control. \textbf{MSS} is the maximum segment size (the payload), distinct
from MTU.

\subsubsection{Segment structure}
Key fields: source/dest ports; 32-bit \textbf{sequence number}
(byte-stream number of the first data byte --- counting \emph{bytes}, not
segments); 32-bit \textbf{acknowledgement number} (next byte expected,
cumulative); header length; flags (\textbf{ACK}, \textbf{RST}/\textbf{SYN}/
\textbf{FIN} for connection management, \textbf{C}/\textbf{E} for congestion
notification); \textbf{receive window} (\texttt{rwnd}, for flow control);
checksum; options.

\subsubsection{Sequence numbers and ACKs}
Telnet example: A types `C'. A $\to$ B: \texttt{Seq=42, ACK=79, data='C'}.
B echoes: \texttt{Seq=79, ACK=43, data='C'}. A acks the echo:
\texttt{Seq=43, ACK=80}. How a receiver handles out-of-order segments is
\emph{not} specified by the standard --- it's up to the implementer.

\subsubsection{RTT estimation and timeout}
The timeout must exceed the RTT, but RTT varies. \texttt{SampleRTT} is
measured per segment (ignoring retransmissions) and smoothed with an
\textbf{exponential weighted moving average (EWMA)}:
\[
    \text{EstimatedRTT} = (1-\alpha)\,\text{EstimatedRTT}
                        + \alpha\,\text{SampleRTT},
    \qquad \alpha = 0.125.
\]
A deviation estimate adds a safety margin:
\[
    \text{DevRTT} = (1-\beta)\,\text{DevRTT}
                  + \beta\,|\text{SampleRTT} - \text{EstimatedRTT}|,
    \qquad \beta = 0.25,
\]
\[
    \boxed{\;\text{TimeoutInterval} = \text{EstimatedRTT} + 4\,\text{DevRTT}\;}.
\]

\subsubsection{Reliable transfer and fast retransmit}
A single timer is kept for the oldest unACKed segment. On a \textbf{timeout},
TCP retransmits the segment that timed out and restarts the timer. A
cumulative ACK can ``cover'' an earlier lost ACK.

\textbf{Fast retransmit:} a missing segment shows up as repeated cumulative
ACKs for the same byte. On receiving \textbf{three duplicate ACKs} (``triple
duplicate''), TCP resends the missing segment \emph{immediately}, without
waiting for the timeout. The receiver's ACK rules (RFC 5681) include
\emph{delayed ACKs} (wait up to 500 ms to piggyback) and immediate
\emph{duplicate ACKs} on out-of-order arrivals.

\subsection{TCP Flow Control}

\emph{Problem:} the network layer may hand data to the receiver faster than
the application reads it from the socket buffer. \textbf{Flow control} stops
the sender from overflowing the \emph{receiver}'s buffer:
\begin{itemize}
    \item The receiver advertises free buffer space in the \texttt{rwnd}
    field of every segment.
    \item The sender limits its in-flight (unACKed) data to $\leq$
    \texttt{rwnd}, guaranteeing the receive buffer never overflows.
\end{itemize}
\textbf{Distinct from congestion control:} flow control protects one
\emph{receiver}; congestion control protects the \emph{network}.

\subsection{TCP Connection Management}

\subsubsection{Why 2-way handshakes fail}
A 2-way handshake can't reliably establish state because of variable delays,
retransmitted/duplicated connection requests, and reordering --- leading to
\textbf{half-open connections} (server thinks it's connected, no client) or
duplicate-data acceptance.

\subsubsection{The 3-way handshake}
\begin{enumerate}
    \item Client $\to$ server: \texttt{SYN}, \texttt{Seq=x} (choose initial
    seq $x$).
    \item Server $\to$ client: \texttt{SYNACK}, \texttt{Seq=y, ACK=x+1}
    (choose initial seq $y$, ack the SYN).
    \item Client $\to$ server: \texttt{ACK, ACK=y+1} (may carry data).
\end{enumerate}
States transition \texttt{LISTEN} $\to$ \texttt{SYN\_SENT}/\texttt{SYN\_RCVD}
$\to$ \texttt{ESTAB}. (``On belay? / Belay on. / Climbing.'')

\begin{lstlisting}[language=Python, caption={The socket calls that drive the TCP handshake}]
# Server
import socket
s = socket.socket(socket.AF_INET, socket.SOCK_STREAM)
s.bind(('', 12000))
s.listen(1)
conn, addr = s.accept()        # completes the 3-way handshake

# Client
c = socket.socket(socket.AF_INET, socket.SOCK_STREAM)
c.connect(('serverName', 12000))   # sends SYN, receives SYNACK, sends ACK
\end{lstlisting}

\subsubsection{Closing a connection}
Each side closes its own half by sending a segment with \texttt{FIN=1} and
responding to a received FIN with an ACK (which may be combined with its own
FIN). Simultaneous FIN exchanges are handled.

\subsection{Principles of Congestion Control}

\textbf{Congestion:} ``too many sources sending too much data too fast for
the network to handle.'' Manifestations: long queueing delays and packet loss
(buffer overflow). A top-10 networking problem --- and \emph{different} from
flow control.

\begin{itemize}
    \item \textbf{Scenario 1} (two flows, one router, infinite buffers): as
    arrival rate $\to R/2$, delay $\to \infty$. Max per-connection throughput
    is $R/2$.
    \item \textbf{Scenario 2} (finite buffers, retransmissions): the
    transport-layer input $\lambda'_{in} \geq \lambda_{in}$ because of
    retransmissions. Costs of congestion: more work per delivered packet, and
    \emph{un-needed duplicates} (from premature timeouts) waste capacity and
    lower achievable throughput.
    \item \textbf{Scenario 3} (multi-hop): when a packet is dropped
    \emph{downstream}, all the upstream capacity and buffering already spent
    on it is wasted.
\end{itemize}

\noindent\textbf{Two approaches:}
\begin{itemize}
    \item \textbf{End-to-end} (TCP): no explicit network feedback; congestion
    \emph{inferred} from observed loss/delay.
    \item \textbf{Network-assisted:} routers give direct feedback
    (e.g.\ TCP ECN, ATM, DECbit), indicating congestion level or an explicit
    rate.
\end{itemize}

\subsection{TCP Congestion Control}

\subsubsection{AIMD}
Senders probe for bandwidth: \textbf{Additive Increase} (raise \texttt{cwnd}
by 1 MSS every RTT while no loss) and \textbf{Multiplicative Decrease}
(halve \texttt{cwnd} on a loss event) --- the characteristic ``sawtooth.''
\begin{itemize}
    \item \textbf{TCP Reno:} on triple-duplicate-ACK loss, cut to half.
    \item \textbf{TCP Tahoe:} on a \emph{timeout}, cut all the way to 1 MSS.
\end{itemize}
AIMD is a distributed, asynchronous algorithm with good stability and
network-wide fairness properties.

\subsubsection{Sending rate and the window}
The sender limits unACKed data to the congestion window:
\[
    \text{LastByteSent} - \text{LastByteAcked} \;\leq\; \texttt{cwnd},
    \qquad \text{TCP rate} \approx \frac{\texttt{cwnd}}{RTT}\ \text{bytes/sec.}
\]

\subsubsection{Slow start, congestion avoidance, fast recovery}
\begin{itemize}
    \item \textbf{Slow start:} begin with \texttt{cwnd}$=1$ MSS and
    \emph{double} \texttt{cwnd} every RTT (increment per ACK) ---
    exponential ramp.
    \item \textbf{Switch to congestion avoidance} when \texttt{cwnd}
    reaches \texttt{ssthresh}; thereafter increase \emph{linearly}
    ($+1$ MSS per RTT). On any loss, \texttt{ssthresh} is set to half the
    \texttt{cwnd} at the loss event.
    \item \textbf{Loss handling:} timeout $\to$ \texttt{ssthresh}$=$
    \texttt{cwnd}$/2$, \texttt{cwnd}$=1$ MSS (back to slow start). Triple
    duplicate ACK $\to$ \texttt{ssthresh}$=$\texttt{cwnd}$/2$,
    \texttt{cwnd}$=$\texttt{ssthresh}$+3$ (\textbf{fast recovery}).
\end{itemize}

\subsubsection{TCP throughput}
Ignoring slow start, if loss occurs at window $W$, the window oscillates
between $W/2$ and $W$, so the average window is $\tfrac{3}{4}W$:
\[
    \text{avg.\ throughput} \approx \frac{3}{4}\cdot\frac{W}{RTT}.
\]
The \textbf{Mathis} formula relates throughput to loss probability $L$:
\[
    \text{throughput} \approx \frac{1.22 \cdot \text{MSS}}{RTT\,\sqrt{L}}.
\]
For ``long, fat pipes'' (e.g.\ 1500-byte segments, $RTT=100$ ms, target
10 Gbps) you'd need $\sim$83\,333 in-flight segments and a tiny loss rate
$L \approx 2\times 10^{-10}$ --- motivating high-speed TCP variants.

\subsubsection{TCP CUBIC}
Instead of pure linear increase, after halving the window on loss, ramp back
toward $W_{\max}$ (the window at which loss was last seen)
\emph{quickly at first, then cautiously} as you near $W_{\max}$. Window growth
is a \emph{cubic} function of the time since the last loss ($K$ is when
\texttt{cwnd} would reach $W_{\max}$). CUBIC is the Linux default and the most
popular TCP for busy web servers.

\subsubsection{Delay-based congestion control (and BBR)}
Goal: ``keep the end-to-end pipe \emph{just} full, but not fuller'' --- keep
the bottleneck busy without inducing loss/large queues. Track
$RTT_{\min}$ (uncongested) and compute uncongested throughput
$=$ \texttt{cwnd}$/RTT_{\min}$:
\begin{itemize}
    \item measured throughput close to uncongested $\Rightarrow$ increase
    \texttt{cwnd} linearly (path not congested);
    \item measured throughput far below $\Rightarrow$ decrease \texttt{cwnd}
    (path congested).
\end{itemize}
This achieves high throughput with low delay and \emph{without forcing loss}.
\textbf{BBR} (a delay/bandwidth-based design) is deployed on Google's
backbone.

\subsubsection{ECN (Explicit Congestion Notification)}
Network-assisted: a router marks two bits in the IP header (ToS field) to
signal congestion; the destination sets the \texttt{ECE} bit on its ACK to
notify the sender. Involves both IP (marking) and TCP (C/E header bits).

\subsubsection{TCP fairness}
If $K$ sessions share a bottleneck of rate $R$, each ``should'' get $R/K$.
TCP is fair \emph{under idealized assumptions} (same RTT, fixed number of
sessions, all in congestion avoidance): additive increase moves both flows
along a $45^\circ$ line, multiplicative decrease pulls them toward equal
share. Caveats:
\begin{itemize}
    \item \textbf{UDP} apps (some multimedia) ignore congestion control ---
    there's no ``Internet police.''
    \item \textbf{Parallel connections:} an app opening many TCP connections
    grabs a larger share (on a link of rate $R$ with 9 existing connections, a
    new app asking for 1 TCP gets $R/10$, but asking for 11 gets $R/2$).
\end{itemize}

% ==================================================================
\section{The Network Layer: Data Plane}
% ==================================================================

\subsection{Overview: Data Plane vs.\ Control Plane}

The network layer transports a segment from the sending to the receiving
host: the sender \textbf{encapsulates} segments into \textbf{datagrams} and
passes them to the link layer; the receiver delivers segments up to the
transport layer. Crucially, the network layer runs in \emph{every} Internet
device --- hosts \emph{and} routers. A router examines header fields of all
IP datagrams passing through it and moves each from an input port to the
appropriate output port.

\subsubsection{Two key functions}
\begin{itemize}
    \item \textbf{Forwarding} (\emph{local}, per-router): move a packet from a
    router's input link to the correct output link. Analogy: getting through
    a \emph{single} highway interchange.
    \item \textbf{Routing} (\emph{global}): determine the end-to-end path
    packets take from source to destination, via routing algorithms. Analogy:
    planning the whole trip.
\end{itemize}

\subsubsection{The two planes}
\begin{itemize}
    \item \textbf{Data plane:} local, per-router; decides how a datagram
    \emph{arriving} on an input port is forwarded to an output port. Runs in
    \emph{hardware} at nanosecond timescales.
    \item \textbf{Control plane:} network-wide logic; decides how a datagram
    is \emph{routed} end-to-end. Runs in \emph{software} at millisecond
    timescales. Two approaches:
    \begin{itemize}
        \item \textbf{Traditional (per-router):} a routing-algorithm component
        in \emph{every} router interacts to build forwarding tables.
        \item \textbf{SDN (Software-Defined Networking):} a \emph{remote}
        controller computes and installs forwarding tables into all routers.
    \end{itemize}
\end{itemize}

\subsection{Network-Layer Service Models}

The design question: \emph{what service model should the ``channel'' that
carries datagrams offer?} Possible per-datagram guarantees include guaranteed
delivery or guaranteed delivery within a delay bound (e.g.\ $<40$ ms).
Per-\emph{flow} guarantees include in-order delivery, a minimum bandwidth, or
bounded jitter (limits on inter-packet-spacing changes).

\renewcommand{\arraystretch}{1.3}
\noindent\begin{tabularx}{\textwidth}{|l|l|X|X|X|X|}
\hline
\textbf{Architecture} & \textbf{Model} & \textbf{Bandwidth} & \textbf{Loss} & \textbf{Order} & \textbf{Timing} \\
\hline
Internet & best effort & none & no & no & no \\
\hline
ATM & CBR & constant rate & yes & yes & yes \\
\hline
ATM & ABR & guaranteed min & no & yes & no \\
\hline
Internet & Intserv Guar.\ (RFC 1633) & yes & yes & yes & yes \\
\hline
Internet & Diffserv (RFC 2475) & possible & possibly & possibly & no \\
\hline
\end{tabularx}

\medskip
\noindent The Internet's \textbf{best-effort} model makes \emph{no}
guarantees on delivery, timing/order, or bandwidth --- yet it succeeded
because its simplicity enabled wide deployment, sufficient bandwidth
provisioning makes real-time apps ``good enough most of the time,''
replicated CDN/datacenter services sit close to clients, and congestion
control of elastic traffic helps. ``It's hard to argue with success.''

\subsection{Inside a Router}

A generic router has \textbf{input ports}, a high-speed \textbf{switching
fabric}, \textbf{output ports}, and a \textbf{routing processor}. The
forwarding data plane (hardware, nanoseconds) is separate from the
routing/management control plane (software, milliseconds). Analogy: a
roundabout (fabric) with entry stations (input ports), exit roads (output
ports), and a station manager (routing processor).

\subsubsection{Input port functions}
\begin{enumerate}
    \item \textbf{Line termination} (physical layer): bit-level reception.
    \item \textbf{Link-layer protocol} (e.g.\ Ethernet) receive.
    \item \textbf{Lookup, forwarding, queueing:} \emph{decentralized
    switching} --- using header fields, look up the output port in a
    forwarding table held in \emph{input-port} memory (``match plus
    action''), ideally at \textbf{line speed}. If datagrams arrive faster
    than the fabric can absorb them, \textbf{input queueing} occurs.
\end{enumerate}
\textbf{Destination-based} forwarding uses only the destination IP address
(traditional); \textbf{generalized} forwarding uses any set of header fields
(the SDN ``match+action'' model).

\subsection{Longest Prefix Matching (LPM)}

When looking up a destination, a router uses the \textbf{longest address
prefix} that matches. Example forwarding table (\texttt{*} = don't-care):

\renewcommand{\arraystretch}{1.3}
\noindent\begin{tabularx}{\textwidth}{|l|X|c|}
\hline
\textbf{Prefix} & \textbf{Meaning} & \textbf{Interface} \\
\hline
\texttt{11001000 00010111 00010***} & \dots/21 & 0 \\
\hline
\texttt{11001000 00010111 00011000} & \dots/24 & 1 \\
\hline
\texttt{11001000 00010111 00011***} & \dots/21 & 2 \\
\hline
otherwise & default & 3 \\
\hline
\end{tabularx}

\medskip
\noindent A datagram to \texttt{11001000 00010111 00011000 10101010} matches
both the \dots/21 (interface 2) \emph{and} the \dots/24 (interface 1) entry;
LPM picks the \textbf{more specific} /24 $\Rightarrow$ \textbf{interface 1}.
A datagram to \texttt{\dots 00010110 \dots} matches only the first prefix
$\Rightarrow$ interface 0.

LPM is typically done in hardware with \textbf{TCAM} (Ternary
Content-Addressable Memory): present an address and it returns the matching
entry in \emph{one clock cycle} regardless of table size (Cisco Catalyst:
$\sim$1M entries).

\subsection{Switching Fabrics}

The fabric transfers a packet from an input link to an output link. The
\textbf{switching rate} (often stated as a multiple of the line rate) should
ideally be $N$ times the line rate $R$ for $N$ ports (aggregate $NR$). Three
classic fabric types:
\begin{itemize}
    \item \textbf{Via memory (1st generation):} a traditional computer copies
    the packet into system memory under CPU control --- \emph{two} bus
    crossings per datagram; speed limited by memory bandwidth.
    \item \textbf{Via a bus:} the datagram crosses a shared bus directly;
    speed limited by \textbf{bus contention} (e.g.\ a 32 Gbps bus suffices for
    access routers).
    \item \textbf{Via an interconnection network:} e.g.\ a \textbf{crossbar}
    (originally built to connect multiprocessors); multiple packets can
    traverse simultaneously, overcoming the single-bus bottleneck.
\end{itemize}

\subsection{Input Queueing and HOL Blocking}

If the fabric is slower than the combined input rate, packets queue at the
inputs (delay and loss from input-buffer overflow). Worse,
\textbf{output-port contention} means only one packet destined to a given
output can cross at a time; a blocked packet at the \emph{front} of an input
queue stalls everything behind it --- even packets bound for idle outputs.
This is \textbf{Head-of-Line (HOL) blocking}.

\subsection{Output Queueing and Buffering}

Buffering is required when datagrams arrive from the fabric faster than the
output link transmits (rate $R$); overflow causes queueing delay and loss.
Two policy questions arise: a \textbf{drop policy} (which datagram to drop
when full) and a \textbf{scheduling discipline} (which queued datagram to
send next).

\subsubsection{How much buffering?}
\begin{itemize}
    \item \textbf{RFC 3439 rule of thumb:} average buffering $\approx$ typical
    RTT ($\sim$250 ms) $\times$ link capacity $C$
    (e.g.\ 10 Gbps $\Rightarrow$ 2.5 Gbit buffer).
    \item \textbf{More recent (with $N$ flows):} buffering
    $\approx \dfrac{\text{RTT}\times C}{\sqrt{N}}$.
    \item \textbf{Bufferbloat:} \emph{too much} buffering (common in home
    routers) inflates delay --- long RTTs give poor real-time performance and
    sluggish TCP response. Recall delay-based congestion control: ``keep the
    bottleneck just full enough, but no fuller.''
\end{itemize}

\subsection{Buffer Management and Packet Scheduling}

\subsubsection{Buffer management}
\begin{itemize}
    \item \textbf{Drop:} which packet to drop when the buffer is full ---
    \emph{tail drop} (drop the arriving packet) or \emph{priority} drop
    (drop/remove by priority).
    \item \textbf{Marking:} which packets to \emph{mark} to signal congestion
    --- ECN, RED (Random Early Detection).
\end{itemize}

\subsubsection{Scheduling policies}
\begin{itemize}
    \item \textbf{FCFS / FIFO:} transmit in order of arrival.
    \item \textbf{Priority:} classify arrivals into priority classes (using
    any header fields); always serve the highest-priority non-empty queue,
    FCFS within a class.
    \item \textbf{Round Robin (RR):} cyclically scan class queues, sending one
    packet from each in turn.
    \item \textbf{Weighted Fair Queueing (WFQ):} generalized RR; each class
    $i$ has weight $w_i$ and receives a share
    $\dfrac{w_i}{\sum_j w_j}$ of the link each cycle --- giving a per-class
    \textbf{minimum bandwidth guarantee}. (Scheduling choices raise
    net-neutrality questions: who gets the best performance?)
\end{itemize}

\subsection{IP: The Internet Protocol --- Datagram Format}

The host/router network layer comprises the \textbf{IP protocol}
(datagram format, addressing, packet-handling conventions), the
\textbf{ICMP protocol} (error reporting, router signaling), and forwarding
tables built by routing protocols (OSPF, BGP) or an SDN controller.

Key IPv4 header fields:
\begin{itemize}
    \item \textbf{Version}, \textbf{header length}, \textbf{type of service}
    (Diffserv bits 0--5, ECN bits 6--7).
    \item \textbf{Total length} (bytes; max 64 KB, typically $\leq$ 1500).
    \item \textbf{16-bit identifier}, \textbf{flags}, \textbf{fragment
    offset} --- for fragmentation/reassembly.
    \item \textbf{TTL} (remaining hops, decremented at each router).
    \item \textbf{Upper-layer protocol} (e.g.\ TCP, UDP --- used for demux at
    the destination).
    \item \textbf{Header checksum}; \textbf{source} and \textbf{destination}
    32-bit IP addresses; \textbf{options} (rarely used).
\end{itemize}
\textbf{Overhead:} 20 bytes IP $+$ 20 bytes TCP $=$ 40 bytes per segment
(before application data).

\subsection{IP Fragmentation and Reassembly}

Every link has an \textbf{MTU} (Maximum Transmission Unit) --- the largest
link-level frame it can carry --- and different link types have different
MTUs. A too-large IP datagram is \textbf{fragmented} within the network into
several datagrams, and \textbf{reassembled only at the destination}. Header
fields coordinate this:
\begin{itemize}
    \item \textbf{Identifier} --- same value on all fragments of one original
    datagram.
    \item \textbf{Flag (MF, ``more fragments'')} --- 1 on all but the last
    fragment.
    \item \textbf{Fragment offset} --- position of this fragment's data in the
    original, \emph{counted in units of 8 bytes}.
\end{itemize}

\subsubsection{Worked example}
A 4000-byte datagram (20-byte header, 3980 bytes of data) over a link with
MTU $=1500$: each fragment carries at most $\lfloor(1500-20)/8\rfloor\times 8
= 1480$ data bytes.

\renewcommand{\arraystretch}{1.3}
\noindent\begin{tabularx}{\textwidth}{|c|c|c|c|c|}
\hline
\textbf{Fragment} & \textbf{length} & \textbf{ID} & \textbf{MF flag} & \textbf{offset} \\
\hline
1 & 1500 & x & 1 & 0 \\
\hline
2 & 1500 & x & 1 & 185 \quad($1480/8$) \\
\hline
3 & 1040 & x & 0 & 370 \quad($2960/8$) \\
\hline
\end{tabularx}

\begin{lstlisting}[language=Python, caption={Computing IPv4 fragments (offset in 8-byte units)}]
def fragment(total_len, mtu, header=20):
    data_left = total_len - header
    max_data  = ((mtu - header) // 8) * 8      # payload/frag, multiple of 8
    frags, offset_bytes = [], 0
    while data_left > 0:
        take = min(max_data, data_left)
        more = 1 if data_left - take > 0 else 0
        frags.append({'length': take + header,
                      'offset': offset_bytes // 8,   # counted in 8-byte units
                      'MF': more})
        offset_bytes += take
        data_left    -= take
    return frags

for f in fragment(4000, 1500):
    print(f)
# {'length': 1500, 'offset': 0,   'MF': 1}
# {'length': 1500, 'offset': 185, 'MF': 1}
# {'length': 1040, 'offset': 370, 'MF': 0}
\end{lstlisting}

\subsection{IP Addressing, Subnets, and CIDR}

An \textbf{IP address} is a 32-bit identifier for an \textbf{interface} (the
connection between a host/router and a physical link). Routers usually have
several interfaces; a host usually has one or two (wired Ethernet, WiFi).
Addresses are written in \textbf{dotted-decimal} notation
(e.g.\ \texttt{223.1.1.1} $=$ \texttt{11011111 00000001 00000001 00000001}).

\subsubsection{Subnets}
A \textbf{subnet} is a set of interfaces that can reach each other
\emph{without} passing through an intervening router. Addresses have
structure: a high-order \textbf{subnet part} (common to all devices in the
subnet) and a low-order \textbf{host part}. \emph{Recipe:} detach each
interface from its host/router; each resulting isolated ``island'' is a
subnet.

\subsubsection{CIDR}
\textbf{CIDR} (Classless Inter-Domain Routing, ``cider'') lets the subnet
portion be \emph{any} length: \texttt{a.b.c.d/x}, where $x$ is the number of
subnet-part bits (the \textbf{subnet mask}). A \texttt{/24} means the first
24 bits are the subnet part.

\begin{lstlisting}[language=Python, caption={Subnet math and longest-prefix match with the ipaddress module}]
import ipaddress

net = ipaddress.ip_network('200.23.16.0/23')
print(net.network_address)     # 200.23.16.0
print(net.broadcast_address)   # 200.23.17.255
print(net.netmask)             # 255.255.254.0
print(net.num_addresses)       # 512   = 2**(32-23)
print(net.num_addresses - 2)   # 510   usable hosts (minus network+broadcast)

# Longest-prefix match: the most specific matching route wins
routes = [ipaddress.ip_network(c) for c in ('200.23.16.0/20',
                                             '200.23.18.0/23')]
dst = ipaddress.ip_address('200.23.18.55')
best = max((r for r in routes if dst in r), key=lambda r: r.prefixlen)
print(best)                    # 200.23.18.0/23  (the /23 beats the /20)
\end{lstlisting}

\subsection{Getting an Address: DHCP}

Two sub-questions: how a \emph{host} gets its host part, and how a
\emph{network} gets its subnet part.

\subsubsection{Host part --- DHCP}
A host's address is either hard-coded by a sysadmin or, more commonly,
obtained dynamically via \textbf{DHCP} (Dynamic Host Configuration Protocol)
--- ``plug and play.'' DHCP lets a host obtain an address when it joins a
network, renew a lease, reuse freed addresses, and support mobile users. The
four-message exchange (\textbf{DORA}), all broadcast, is:
\begin{enumerate}
    \item \textbf{Discover} --- ``is there a DHCP server out there?''
    (src \texttt{0.0.0.0:68}, dst \texttt{255.255.255.255:67}).
    \item \textbf{Offer} --- server proposes an address (\texttt{yiaddr}) and
    lease time.
    \item \textbf{Request} --- client asks for the offered address.
    \item \textbf{ACK} --- server confirms.
\end{enumerate}
The first two steps can be skipped if a client is reusing a previously
allocated address (RFC 2131). DHCP returns more than an address: the
\textbf{first-hop (default) router}, the \textbf{DNS server}, and the
\textbf{subnet mask}. The whole DHCP message is carried in
UDP\,$\to$\,IP\,$\to$\,Ethernet and broadcast on the LAN.

\subsubsection{Network part --- ISPs and ICANN}
A network gets its subnet part as an allocated block of its provider ISP's
address space. \textbf{ICANN} (via 5 Regional Registries) allocates IP blocks
and manages the DNS root zone and TLD delegation. The last IPv4 blocks were
handed to the registries in 2011 --- motivating NAT and IPv6.

\subsection{Hierarchical Addressing and Route Aggregation}

Hierarchical allocation lets an ISP advertise a \emph{single} aggregate route
(``send me anything beginning \texttt{200.23.16.0/20}'') covering many
customer organizations --- keeping global routing tables small. If one
organization moves to a different ISP, that ISP advertises a \textbf{more
specific} route (e.g.\ \texttt{200.23.18.0/23}); LPM ensures the more specific
route wins, so traffic still reaches the moved organization.

\subsection{NAT: Network Address Translation}

All devices in a local network share \emph{one} public IPv4 address as far as
the outside world is concerned. Internally they use a \textbf{private} block
(\texttt{10.0.0.0/8}, \texttt{172.16.0.0/12}, or \texttt{192.168.0.0/16}) that
is only meaningful locally.

\textbf{Advantages:} only one ISP-provided address is needed; local addresses
can change without notifying the outside; the ISP can change without
renumbering local hosts; and inside devices aren't directly addressable from
outside (a security side-effect).

\subsubsection{How it works}
The NAT router transparently rewrites addresses using a \textbf{translation
table}:
\begin{itemize}
    \item \textbf{Outgoing:} replace (source IP, port) with
    (NAT IP, \emph{new} port), and record the mapping.
    \item \textbf{Incoming:} replace the destination (NAT IP, new port) with
    the stored (source IP, port).
\end{itemize}

\renewcommand{\arraystretch}{1.3}
\noindent\begin{tabularx}{\textwidth}{|X|X|}
\hline
\textbf{WAN side (NAT IP, port)} & \textbf{LAN side (host IP, port)} \\
\hline
\texttt{138.76.29.7, 5001} & \texttt{10.0.0.1, 3345} \\
\hline
\end{tabularx}

\medskip
\noindent \textbf{Controversy:} NAT makes a network-layer device manipulate
port numbers (a transport-layer concept), violating the end-to-end argument,
and complicates \textbf{NAT traversal} (reaching a server \emph{behind} a
NAT). Purists argue address shortage ``should'' be solved by IPv6 --- but NAT
is here to stay (home, institutional, and 4G/5G networks).

\subsection{IPv6}

\textbf{Motivation:} the 32-bit IPv4 space was being exhausted; secondary
goals were faster processing (a fixed 40-byte header) and better handling of
\emph{flows}.

\subsubsection{Datagram format}
Fields: version, \textbf{priority}, \textbf{flow label} (identifies datagrams
in the same flow --- concept loosely defined), payload length, \textbf{next
header} (upper-layer protocol), \textbf{hop limit}, and \textbf{128-bit}
source/destination addresses. Compared with IPv4, IPv6 \emph{drops}:
\begin{itemize}
    \item the header \textbf{checksum} (speeds router processing),
    \item in-network \textbf{fragmentation/reassembly},
    \item \textbf{options} (moved to an upper-layer ``next-header'' chain).
\end{itemize}

\subsubsection{Transition: tunneling}
Not all routers can be upgraded at once (no ``flag days''), so IPv4 and IPv6
must coexist. \textbf{Tunneling} carries an IPv6 datagram as the
\emph{payload} of an IPv4 datagram between IPv6-capable routers separated by
IPv4-only infrastructure (a ``packet within a packet''). The entry router
(dual-stack IPv6/v4) encapsulates the IPv6 datagram in an IPv4 datagram
addressed to the far tunnel end; the exit router decapsulates and forwards the
original IPv6 datagram. The inner IPv6 source/destination
(e.g.\ A$\to$F) are preserved; the outer IPv4 header uses the tunnel
endpoints (e.g.\ B$\to$E). Tunneling is used widely elsewhere too (4G/5G).

\medskip
\noindent \textbf{Adoption} is slow (25+ years): $\sim$40\% of Google's
clients reach it over IPv6 (2023), and $\sim$1/3 of US government domains are
IPv6-capable.

% ==================================================================
\section{The Network Layer: Control Plane}
% ==================================================================

\subsection{Overview: Two Approaches}

The control plane is the network-wide logic that computes the forwarding
tables the data plane uses. Recall the split: \textbf{forwarding} moves a
packet from a router's input to the right output (data plane, local, per
router); \textbf{routing} determines the end-to-end route packets take
(control plane). Two ways to structure it:
\begin{itemize}
    \item \textbf{Per-router control (traditional):} a routing-algorithm
    component runs in \emph{every} router; the components interact (exchange
    messages) and each router computes its own forwarding table.
    \item \textbf{Logically centralized control (SDN):} a remote controller
    computes forwarding tables and installs them into the routers, which
    become simple match+action devices.
\end{itemize}
Internet instantiations: OSPF and BGP (per-router); OpenFlow with ODL/ONOS
controllers (SDN); ICMP for control messaging; SNMP and NETCONF/YANG for
network management.

\subsection{Routing Protocols and the Graph Abstraction}

Goal: find ``good'' (least-cost / fastest / least-congested) paths from source
to destination hosts through the router network --- a top-10 networking
problem. Model the network as a graph $G=(N,E)$:
\begin{itemize}
    \item $N$ = set of routers, e.g.\ $\{u,v,w,x,y,z\}$; $E$ = set of links.
    \item $c_{a,b}$ = cost of the direct link $(a,b)$; $c_{a,b}=\infty$ if
    $a,b$ are not direct neighbors. The operator defines cost: always $1$ (hop
    count), or inversely related to bandwidth or to congestion.
    \item Path cost = sum of link costs; we want the \textbf{least-cost path}.
\end{itemize}
\textbf{Classification.} \emph{Global} information (link-state): every router
knows the complete topology and all link costs. \emph{Decentralized}
(distance-vector): an iterative process where each router initially knows only
the costs to its attached neighbors and exchanges info with them. Orthogonally,
routes are \emph{static} (change slowly) or \emph{dynamic} (periodic updates /
react to cost changes).

\subsection{Link-State Routing: Dijkstra's Algorithm}

Every node floods its link costs network-wide (a \textbf{link-state
broadcast}), so all nodes learn the identical full topology. Dijkstra computes
least-cost paths from one \emph{source} to all other nodes (yielding that
source's forwarding table). It is iterative: after $k$ iterations, the least
costs to $k$ destinations are final. Notation: $D(v)$ = current least-cost
estimate to $v$; $p(v)$ = predecessor of $v$ on that path; $N'$ = set of
finalized nodes.

\begin{lstlisting}[caption={Dijkstra's link-state algorithm (source $u$)}]
Initialization:
  N' = {u}
  for all nodes v:
     if v adjacent to u:  D(v) = c(u,v)     # direct-neighbor cost only
     else:                D(v) = infinity

Loop until all nodes in N':
  find w not in N' with minimum D(w); add w to N'
  update D(v) for all v adjacent to w and not in N':
     D(v) = min( D(v),  D(w) + c(w,v) )
\end{lstlisting}

\subsubsection{Worked example (source $u$)}
Costs: $c_{u,v}{=}2$, $c_{u,x}{=}1$, $c_{u,w}{=}5$, $c_{v,x}{=}2$,
$c_{v,w}{=}3$, \\ $c_{x,w}{=}3$, $c_{x,y}{=}1$, $c_{w,y}{=}1$, $c_{w,z}{=}5$,
$c_{y,z}{=}2$.

\renewcommand{\arraystretch}{1.3}
\noindent\begin{tabularx}{\textwidth}{|c|l|c|c|c|c|c|}
\hline
\textbf{Step} & \textbf{$N'$} & \textbf{$D(v),p$} & \textbf{$D(w),p$} & \textbf{$D(x),p$} & \textbf{$D(y),p$} & \textbf{$D(z),p$} \\
\hline
0 & $u$      & 2,$u$ & 5,$u$ & \textbf{1,$u$} & $\infty$ & $\infty$ \\
\hline
1 & $ux$     & \textbf{2,$u$} & 4,$x$ & & 2,$x$ & $\infty$ \\
\hline
2 & $uxy$    & 2,$u$ & 3,$y$ & & & 4,$y$ \\
\hline
3 & $uxyv$   & & \textbf{3,$y$} & & & 4,$y$ \\
\hline
4 & $uxyvw$  & & & & & \textbf{4,$y$} \\
\hline
5 & $uxyvwz$ & & & & & \\
\hline
\end{tabularx}

\medskip
\noindent Resulting least costs from $u$: $x{=}1$, $v{=}2$, $y{=}2$
($u$-$x$-$y$), $w{=}3$ ($u$-$x$-$y$-$w$), $z{=}4$ ($u$-$x$-$y$-$z$). The
forwarding table sends $v$ out link $(u,v)$ and everything else out $(u,x)$.

\begin{lstlisting}[language=Python, caption={Dijkstra's least-cost paths}]
import heapq

def dijkstra(graph, src):          # graph[u] = {v: cost, ...}
    dist = {n: float('inf') for n in graph}
    prev = {n: None for n in graph}
    dist[src] = 0
    pq = [(0, src)]
    while pq:
        d, u = heapq.heappop(pq)
        if d > dist[u]:            # stale entry
            continue
        for v, w in graph[u].items():
            if d + w < dist[v]:
                dist[v], prev[v] = d + w, u
                heapq.heappush(pq, (dist[v], v))
    return dist, prev
\end{lstlisting}

\noindent \textbf{Complexity.} $n$ iterations, each checking the remaining
nodes $\Rightarrow n(n{+}1)/2$ comparisons $= O(n^2)$; $O(n\log n)$ with a
heap. \textbf{Message complexity:} each router broadcasts its link state
(each broadcast crosses $O(n)$ links) $\Rightarrow O(n^2)$ overall.
\textbf{Oscillations} are possible when link costs depend on carried traffic
(routes ``slosh'' between paths).

\subsection{Distance-Vector Routing: Bellman-Ford}

Based on the \textbf{Bellman-Ford equation} (dynamic programming). Let $D_x(y)$
be the least cost from $x$ to $y$:
\[
    D_x(y) = \min_{v}\big\{\, c_{x,v} + D_v(y) \,\big\},
\]
the minimum over all neighbors $v$ of ($x$-to-$v$ direct cost) $+$
($v$'s least cost to $y$). The neighbor $v$ achieving the minimum is the
\textbf{next hop}.

\subsubsection{Worked example}
Reaching $z$ from $u$ with neighbors $v,x,w$, where $D_v(z){=}5$, $D_x(z){=}3$,
$D_w(z){=}3$ and $c_{u,v}{=}2$, $c_{u,x}{=}1$, $c_{u,w}{=}5$:
\[
    D_u(z) = \min\{\,2{+}5,\ 1{+}3,\ 5{+}3\,\} = \min\{7,4,8\} = 4,
\]
so the next hop toward $z$ is $x$.

\textbf{The algorithm is iterative, asynchronous, distributed, and
self-stopping.} Each node keeps its own distance vector plus its neighbors'
vectors. A local iteration is triggered by a \emph{local link-cost change} or a
\emph{DV update from a neighbor}; the node recomputes with Bellman-Ford, and
\emph{only if its DV changed} does it notify its neighbors --- who then repeat,
only if necessary. No message, no action. Under natural conditions the
estimates $D_x(y)$ converge to the true least costs $d_x(y)$.

\begin{lstlisting}[language=Python, caption={A distance-vector update at node x (one Bellman-Ford step)}]
def dv_update(x, link_cost, neighbor_dv, destinations):
    # link_cost[v]      = cost of the direct x->v link (v a neighbor)
    # neighbor_dv[v][y] = neighbor v's estimated cost to y
    Dx = {}
    for y in destinations:
        Dx[y] = min(link_cost[v] + neighbor_dv[v][y] for v in link_cost)
    return Dx        # notify neighbors iff Dx changed
\end{lstlisting}

\subsubsection{Convergence example (3 nodes)}
Costs $c_{x,y}{=}2$, $c_{x,z}{=}7$, $c_{y,z}{=}1$. From its neighbors' vectors,
node $x$ computes
\[
    D_x(y)=\min\{c_{x,y}{+}D_y(y),\; c_{x,z}{+}D_z(y)\}=\min\{2{+}0,\,7{+}1\}=2,
\]
\[
    D_x(z)=\min\{c_{x,y}{+}D_y(z),\; c_{x,z}{+}D_z(z)\}=\min\{2{+}1,\,7{+}0\}=3.
\]
After exchanging vectors, all three settle (here $D_x=[0,2,3]$). Information
diffuses one hop per round: a node's state at $t{=}0$ can influence
computations up to $k$ hops away by $t{=}k$.

\subsubsection{Link-cost changes: good news vs.\ bad news}
\begin{itemize}
    \item \textbf{Good news travels fast:} when a link cost \emph{drops}, a
    node updates its DV and informs neighbors; the news settles in a few
    rounds.
    \item \textbf{Bad news travels slowly --- count-to-infinity:} suppose
    $y$'s link to $x$ jumps to $60$ while $z$ still advertises a cost-$5$ path.
    $y$ computes ``$6$ via $z$'' and tells $z$; $z$ then computes ``$7$ via
    $y$'' and tells $y$; $y \to 8$; $z \to 9$; \dots incrementing slowly toward
    the true cost.
    \item \textbf{Poisoned reverse:} if $x$ routes to a destination via $z$,
    $x$ advertises its distance to that destination as $\infty$ to $z$, so $z$
    won't loop back through $x$. Fixes 2-node loops, not larger ones.
\end{itemize}

\renewcommand{\arraystretch}{1.3}
\noindent\begin{tabularx}{\textwidth}{|l|X|X|}
\hline
 & \textbf{Link-State (LS)} & \textbf{Distance-Vector (DV)} \\
\hline
Knowledge & Full topology (flooded) & Neighbors + their vectors \\
\hline
Algorithm & Dijkstra & Bellman-Ford \\
\hline
Messages & $O(n^2)$; each floods its link state & Only between neighbors; varies \\
\hline
Convergence & $O(n^2)$; possible oscillations & Varies; routing loops, count-to-infinity \\
\hline
Robustness & Router can advertise wrong \emph{own} cost; each computes its own table & A router can advertise wrong \emph{path} cost (``low cost to everywhere'' = black-holing); errors \emph{propagate} \\
\hline
\end{tabularx}

\subsection{Making Routing Scalable: Autonomous Systems}

The idealized ``all routers identical, flat network'' does not scale: billions
of destinations can't fit in one table, and exchanging them would swamp links.
Also, each network's admin wants \textbf{autonomy} over its own routing. So
routers are aggregated into \textbf{Autonomous Systems (ASes)} / domains:
\begin{itemize}
    \item \textbf{Intra-AS (intra-domain):} routing among routers in the same
    AS; all its routers run the same intra-AS protocol.
    \item \textbf{Inter-AS (inter-domain):} routing between ASes, done by
    \textbf{gateway routers} at the AS edge (which also run intra-AS routing).
\end{itemize}
Intra-AS routing fills table entries for internal destinations; inter-AS plus
intra-AS together fill entries for external destinations. When a router in AS1
receives a datagram for outside AS1, inter-domain routing must (1) learn which
destinations are reachable through AS2 vs.\ AS3 and (2) propagate that
reachability to all AS1 routers (so a packet is sent to the right gateway).

\subsection{Intra-AS Routing: RIP, EIGRP, OSPF}

\begin{itemize}
    \item \textbf{RIP} (RFC 1723): classic distance-vector, DVs exchanged every
    30 s; no longer widely used.
    \item \textbf{EIGRP:} DV-based, long Cisco-proprietary (opened in 2013,
    RFC 7868).
    \item \textbf{OSPF} (Open Shortest Path First, RFC 2328): the dominant IGP;
    IS-IS (an ISO standard) is essentially equivalent.
\end{itemize}

\textbf{OSPF} is ``open'' (publicly specified) and a classic \emph{link-state}
protocol: each router floods link-state advertisements to all others in the AS
\emph{directly over IP} (not TCP/UDP), builds the full topology, and runs
Dijkstra to compute its forwarding table. Cost metrics are configurable
(bandwidth, delay). All OSPF messages are \textbf{authenticated} to prevent
malicious intrusion.

\textbf{Hierarchical OSPF} uses two levels --- local \emph{areas} and a
\emph{backbone}. LSAs are flooded only within an area (or the backbone); each
node knows its area's detailed topology and only the \emph{direction} to other
areas.
\begin{itemize}
    \item \textbf{Area border routers} summarize their area's distances and
    advertise them into the backbone.
    \item \textbf{Backbone routers} run OSPF limited to the backbone.
    \item \textbf{Boundary routers} connect to other ASes.
\end{itemize}

\subsection{Inter-AS Routing: BGP}

\textbf{BGP} (Border Gateway Protocol) is the de-facto inter-domain routing
protocol --- ``the glue that holds the Internet together.'' It lets a subnet
advertise its existence and the destinations it can reach: ``I am here, here's
who I can reach, and how.'' BGP gives each AS a way to:
\begin{itemize}
    \item obtain destination reachability from neighboring ASes (\textbf{eBGP});
    \item determine routes based on reachability \emph{and policy};
    \item propagate reachability to all AS-internal routers (\textbf{iBGP});
    \item advertise reachability to neighboring networks.
\end{itemize}
Gateway routers run \emph{both} eBGP (to other ASes) and iBGP (inside the AS).

\subsubsection{BGP sessions and messages}
Two BGP peers exchange messages over a semi-permanent \textbf{TCP} connection.
BGP is a \textbf{path-vector} protocol: advertisements carry \emph{paths} to
destination prefixes. Messages (RFC 4271): \texttt{OPEN} (open + authenticate
the peer), \texttt{UPDATE} (advertise a new path or withdraw an old one),
\texttt{KEEPALIVE} (keep the session alive; also ACKs OPEN),
\texttt{NOTIFICATION} (report errors / close).

\subsubsection{Path attributes}
An advertised route $=$ \textbf{prefix} $+$ \textbf{attributes}. Two key
attributes:
\begin{itemize}
    \item \textbf{AS-PATH:} the list of ASes the advertisement has traversed
    (also used to detect and avoid loops).
    \item \textbf{NEXT-HOP:} the specific internal-AS router leading toward the
    next-hop AS.
\end{itemize}
\textbf{Policy-based routing:} an \emph{import} policy accepts/declines a
received path (e.g.\ ``never route through AS $Y$''); an \emph{export} policy
decides whether to advertise a path onward. When AS3's gateway advertises
$\langle AS3,X\rangle$ to AS2, AS3 \emph{promises} to forward datagrams toward
$X$.

\subsubsection{Path propagation and forwarding tables}
AS2 router 2c receives $\langle AS3,X\rangle$ via eBGP from 3a, propagates
$\langle AS2,AS3,X\rangle$ inside AS2 via iBGP, and 2a advertises it to AS1 via
eBGP. AS1's gateway 1c may learn \emph{two} paths to $X$
($\langle AS2,AS3,X\rangle$ and $\langle AS3,X\rangle$), pick one by policy, and
flood it via iBGP. To build the actual forwarding table, a router combines BGP
(``$X$ is reachable via gateway 1c'') with intra-AS OSPF (``to reach 1c, use
this interface''): $X$-bound traffic leaves on the interface OSPF says leads to
that gateway.

\subsubsection{Hot-potato routing and policy}
\textbf{Hot-potato routing:} among multiple gateways that can reach $X$, pick
the one with least \emph{intra-domain} cost --- get the packet out of your own
AS as soon as possible, regardless of inter-domain hop count.
\textbf{Policy example:} an ISP typically routes traffic only to/from its
\emph{customer} networks and refuses to carry \emph{transit} traffic between
other ISPs (no revenue). So a provider $B$ that learns a path may choose
\emph{not} to advertise it if none of the endpoints are $B$'s customers, and a
dual-homed customer $x$ won't advertise a route that would make it transit
between two providers.

\textbf{Route selection order} (when multiple routes to a destination exist):
\begin{enumerate}
    \item highest \textbf{local-preference} (a policy decision);
    \item shortest \textbf{AS-PATH};
    \item closest \textbf{NEXT-HOP} (hot-potato routing);
    \item additional tie-breakers.
\end{enumerate}

\subsubsection{Why intra- and inter-AS routing differ}
\textbf{Policy:} inter-AS admins want control over who routes through them
(policy dominates); intra-AS is a single admin, so policy matters less.
\textbf{Scale:} hierarchy shrinks table size and update traffic.
\textbf{Performance:} intra-AS can optimize for performance; inter-AS is
dominated by policy.

\subsection{The SDN Control Plane}

SDN separates a \emph{logically centralized} controller from simple data-plane
switches:
\begin{itemize}
    \item \textbf{Generalized forwarding:} switches hold \textbf{flow tables}
    of \textbf{match+action} rules (\textbf{OpenFlow}); a rule matches header
    fields across layers 2--4 and acts (forward to port(s), drop, modify a
    field, or send to the controller).
    \item \textbf{Controller:} maintains network state and offers a
    \textbf{southbound} API (e.g.\ OpenFlow) to program switches and a
    \textbf{northbound} API to control apps (routing, access control, load
    balancing). Production controllers include OpenDaylight (ODL) and ONOS.
\end{itemize}
Benefits: programmability, one place for network-wide logic, and open
interfaces (no vendor lock-in).

\subsection{ICMP: Internet Control Message Protocol}

Hosts and routers use \textbf{ICMP} to exchange network-level information ---
error reporting (unreachable host / network / port / protocol) and echo
request/reply (used by \texttt{ping}). ICMP sits ``above'' IP but its messages
are carried \emph{in} IP datagrams; each message has a \textbf{type}, a
\textbf{code}, plus the first 8 bytes of the IP datagram that caused the error.

\renewcommand{\arraystretch}{1.25}
\noindent\begin{tabularx}{\textwidth}{|c|c|X||c|c|X|}
\hline
\textbf{Type} & \textbf{Code} & \textbf{Description} & \textbf{Type} & \textbf{Code} & \textbf{Description} \\
\hline
0 & 0 & echo reply (ping) & 3 & 6 & dest.\ network unknown \\
\hline
3 & 0 & dest.\ network unreachable & 3 & 7 & dest.\ host unknown \\
\hline
3 & 1 & dest.\ host unreachable & 8 & 0 & echo request (ping) \\
\hline
3 & 2 & dest.\ protocol unreachable & 9 & 0 & route advertisement \\
\hline
3 & 3 & dest.\ port unreachable & 10 & 0 & router discovery \\
\hline
4 & 0 & source quench (unused) & 11 & 0 & TTL expired \\
\hline
 & & & 12 & 0 & bad IP header \\
\hline
\end{tabularx}

\subsubsection{Traceroute via ICMP}
The source sends sets of UDP segments with increasing TTL (first set TTL$=1$,
next TTL$=2$, \dots). The datagram in the $n$-th set expires at router $n$,
which discards it and returns an ICMP \textbf{TTL-expired} (type 11, code 0),
often naming the router; the source records the RTT. It \emph{stops} when a
segment finally reaches the destination host, which returns \textbf{port
unreachable} (type 3, code 3).

\subsection{Network Management: SNMP and NETCONF/YANG (Brief)}

A \textbf{managing server} monitors and controls \textbf{managed devices}, each
running an \textbf{agent} that exposes a \textbf{MIB} (Management Information
Base --- the data objects). \textbf{SNMP} carries request/response
(\texttt{GET}/\texttt{SET}) messages and unsolicited \textbf{trap} messages
(agent$\to$server event reports). \textbf{NETCONF} (with the \textbf{YANG}
data-modeling language) is the modern approach to configuring and managing
devices.

% ==================================================================
\section{The Link Layer and LANs}
% ==================================================================

\subsection{Introduction and Services}

Terminology: \textbf{nodes} = hosts and routers; \textbf{links} connect
adjacent nodes (wired, wireless, LANs); a link-layer (Layer-2) packet is a
\textbf{frame} that encapsulates a datagram. The link layer's job is to move a
datagram from one node to a \emph{physically adjacent} node over a single link.
A datagram may traverse different link protocols on different links (WiFi on the
first, Ethernet on the next), each offering different services.
\emph{Transportation analogy:} a trip Princeton$\to$Lausanne (limo, plane,
train) --- the tourist is the datagram, each leg is a link, each mode is a
link-layer protocol, and the travel agent is the routing algorithm.

\textbf{Possible services:}
\begin{itemize}
    \item \textbf{Framing / link access:} encapsulate the datagram (header +
    trailer); arbitrate access if the medium is shared; put source/destination
    \textbf{MAC addresses} (distinct from IP addresses) in the header.
    \item \textbf{Reliable delivery} between adjacent nodes --- seldom used on
    low-error wired links, valuable on high-error wireless links. (Why both
    link-level and end-to-end reliability? To avoid costly end-to-end
    retransmissions across a single lossy link.)
    \item \textbf{Flow control}, \textbf{error detection}, \textbf{error
    correction}, and \textbf{half-/full-duplex} operation.
\end{itemize}
The link layer lives in the \textbf{network adapter (NIC)} --- on-chip or a
card implementing the link and physical layers --- a combination of hardware,
software, and firmware attached to the host's system bus. Sending side:
encapsulate, add error-check bits. Receiving side: check for errors, extract
the datagram, pass it up.

\subsection{Error Detection and Correction}

The sender appends \textbf{EDC} (error-detection-and-correction) bits to the
data $D$. Detection is \emph{not} 100\% reliable (some errors slip through
undetected, though rarely); a larger EDC field improves detection/correction.
\begin{itemize}
    \item \textbf{Single-bit parity:} one bit makes the number of 1s even (or
    odd); detects an \emph{odd} number of bit errors only.
    \item \textbf{Two-dimensional parity:} parity per row \emph{and} per
    column; \emph{detects and corrects} any single-bit error by locating its
    row and column --- forward error correction, no retransmission.
    \item \textbf{Internet checksum:} the one's-complement sum from the
    transport layer --- software, relatively weak.
    \item \textbf{CRC (Cyclic Redundancy Check):} stronger, hardware-friendly,
    used by Ethernet and 802.11 WiFi.
\end{itemize}

\subsubsection{CRC}
Given data $D$ and a standard $(r{+}1)$-bit generator $G$, the sender finds the
$r$-bit remainder
\[
    R = \text{remainder}\!\left(\frac{D\cdot 2^{r}}{G}\right)\ (\text{mod } 2),
\]
so that $D\cdot 2^{r}\oplus R$ is exactly divisible by $G$ (mod 2); it transmits
$D\cdot 2^{r}\oplus R$. The receiver divides the received bits by $G$; a
\emph{nonzero} remainder means an error. CRC detects all burst errors of length
$\leq r$. \textbf{Example:} $D=101110$, $G=1001$ ($r=3$); mod-2 dividing
$101110000$ by $1001$ gives $R=011$, so the transmitted frame is
$101110\,011$.

\begin{lstlisting}[language=Python, caption={Mod-2 (CRC) remainder}]
def crc_remainder(data_bits, gen_bits):
    G = int(gen_bits, 2)
    r = len(gen_bits) - 1
    reg = int(data_bits, 2) << r            # append r zeros: D * 2^r
    for i in range(len(data_bits) + r - 1, r - 1, -1):
        if reg & (1 << i):                  # subtract (XOR) G aligned at bit i
            reg ^= G << (i - r)
    return format(reg, f'0{r}b')

print(crc_remainder('101110', '1001'))      # -> '011'
\end{lstlisting}

\subsection{Multiple Access Protocols}

Two link types: \textbf{point-to-point} (one sender/receiver, e.g.\
switch-to-host, PPP dial-up) and \textbf{broadcast} (a shared medium: old
Ethernet, upstream cable HFC, 802.11 WiFi, 4G/5G, satellite). On a broadcast
link, two or more simultaneous transmissions \textbf{collide} (a node receiving
two signals at once). A multiple-access (MAC) protocol is a \emph{distributed}
algorithm that decides when a node may transmit --- and coordination must use
the shared channel itself (no separate out-of-band channel).

\textbf{Ideal MAC on a rate-$R$ channel:} (1) one active node gets the full
rate $R$; (2) $M$ active nodes each get $\approx R/M$; (3) fully decentralized
(no coordinator, no clock sync); (4) simple. Three families:

\subsubsection{Channel partitioning}
\begin{itemize}
    \item \textbf{TDMA} (time-division): each station gets a fixed slot per
    round (slot $=$ one frame time); unused slots go idle (6-station LAN with
    only 1,3,4 active $\Rightarrow$ slots 2,5,6 wasted).
    \item \textbf{FDMA} (frequency-division): each station gets a fixed
    frequency band; idle bands wasted.
\end{itemize}
Fair and collision-free, but inefficient at low load (a node is capped at
$R/N$ even when it is the only one active).

\subsubsection{Random access}
No coordination; collisions happen and must be recovered from.
\begin{itemize}
    \item \textbf{Slotted ALOHA:} synchronized slots (one frame each); a node
    with a fresh frame transmits in the next slot; on collision it retransmits
    in each later slot with probability $p$ until success. \emph{Pros:} a lone
    node uses the full rate; highly decentralized; simple. \emph{Cons:}
    collisions and idle slots waste slots; needs clock sync.
    \item \textbf{Pure ALOHA:} unsynchronized (simpler, no slots) but more
    collisions.
    \item \textbf{CSMA} (carrier sense): \emph{listen before transmit} --- send
    only if the channel is idle. Collisions still occur because of propagation
    delay (two nodes sense idle and both start); the whole colliding frame time
    is wasted.
    \item \textbf{CSMA/CD} (collision detection): abort as soon as a collision
    is heard, cutting wasted time. Used by classic \textbf{Ethernet}. Detection
    is easy on wire, hard on wireless.
    \item \textbf{CSMA/CA} (collision avoidance): used by WiFi (802.11), where
    a node cannot detect collisions while transmitting.
\end{itemize}

\textbf{Slotted-ALOHA efficiency.} With $N$ nodes each transmitting in a slot
with probability $p$,
\[
    P(\text{a given node succeeds})=p(1-p)^{N-1},\qquad
    P(\text{some node succeeds})=Np(1-p)^{N-1}.
\]
Maximizing over $p$ and letting $N\to\infty$ gives max efficiency
$1/e\approx 0.37$ --- at best, the channel does useful work 37\% of the time
(pure ALOHA: $1/(2e)\approx 0.18$).

\begin{lstlisting}[caption={Ethernet CSMA/CD with binary exponential backoff}]
1. Adapter gets a datagram from the network layer, creates a frame.
2. Sense the channel: if idle, start transmitting; if busy, wait until idle.
3. Whole frame sent with no collision  ->  done.
4. Collision detected while sending     ->  abort, send a jam signal.
5. Binary exponential backoff: after the m-th collision, pick K uniformly
   from {0, 1, ..., 2^m - 1}; wait K * 512 bit-times; go to step 2.
   (More collisions  ->  larger backoff window.)
\end{lstlisting}

\subsubsection{Taking turns}
\begin{itemize}
    \item \textbf{Polling:} a master node invites each node to transmit in turn
    (polling overhead, latency, single point of failure; Bluetooth uses
    polling).
    \item \textbf{Token passing:} a control token is passed node-to-node; only
    its holder may transmit (token overhead/latency, token loss; FDDI, token
    ring).
\end{itemize}
``Taking turns'' aims for the best of both worlds --- efficient at high load
(like partitioning) yet low-delay at low load (like random access).

\subsection{MAC Addressing and ARP}

\subsubsection{MAC addresses}
An interface's \textbf{MAC (LAN / physical / Ethernet) address} is 48 bits
(most LANs), burned into the NIC ROM (sometimes software-settable), written in
hex (e.g.\ \texttt{1A-2F-BB-76-09-AD}). Its job: get a frame from one interface
to another \emph{physically-connected} interface on the same subnet. Allocation
is administered by \textbf{IEEE} (manufacturers buy address blocks for
uniqueness). Analogy: \textbf{MAC $=$ Social-Security number} (flat, portable
--- moves with the adapter), \textbf{IP $=$ postal address} (hierarchical,
depends on the attached subnet, \emph{not} portable).

\subsubsection{ARP}
\textbf{ARP} (Address Resolution Protocol) answers: knowing an interface's IP
address on my LAN, what is its MAC address? Each node keeps an \textbf{ARP
table} of $\langle$IP, MAC, TTL$\rangle$ entries (TTL typically $\sim$20 min).
When $A$ wants to send to $B$ but lacks $B$'s MAC:
\begin{enumerate}
    \item $A$ \textbf{broadcasts} an ARP query (destination MAC
    \texttt{FF-FF-FF-FF-FF-FF}) containing $B$'s IP address; all LAN nodes
    receive it.
    \item $B$ recognizes its own IP and \textbf{replies (unicast)} to $A$ with
    its MAC address.
    \item $A$ caches the $\langle$IP, MAC$\rangle$ mapping (with a TTL) in its
    ARP table.
\end{enumerate}

\subsubsection{Routing to another subnet}
To reach a host $B$ on \emph{another} subnet via first-hop router $R$: $A$ knows
$B$'s IP, $R$'s IP, and (via ARP) $R$'s MAC. The IP source/destination stay
$A$/$B$ end-to-end; only the \emph{link-layer} addresses change hop by hop.

\renewcommand{\arraystretch}{1.3}
\noindent\begin{tabularx}{\textwidth}{|l|l|l|X|}
\hline
\textbf{Hop} & \textbf{IP src $\to$ dst} & \textbf{MAC src $\to$ dst} & \textbf{Comment} \\
\hline
$A \to R$ & $A \to B$ & $\text{MAC}_A \to \text{MAC}_R^{\text{(left)}}$ & Frame carries $A$-to-$B$ datagram; dest MAC is $R$'s \\
\hline
$R \to B$ & $A \to B$ & $\text{MAC}_R^{\text{(right)}} \to \text{MAC}_B$ & $R$ removes/re-adds Layer-2 header; dest MAC is $B$'s \\
\hline
\end{tabularx}

\medskip
\noindent Key insight: routers rewrite the Layer-2 header at each hop, while the
Layer-3 (IP) destination remains the true endpoint $B$ throughout.

\subsection{Ethernet}

The dominant, cheap, simple wired LAN, keeping up with the speed race
(10 Mbps--400 Gbps). The physical topology evolved from a \textbf{bus} (all
nodes in one collision domain, popular through the mid-90s) to today's
\textbf{switched} star: an active Layer-2 switch at the center, each spoke its
own collision domain, so nodes do not collide.

\subsubsection{Frame structure}
\renewcommand{\arraystretch}{1.3}
\noindent\begin{tabularx}{\textwidth}{|l|l|X|}
\hline
\textbf{Field} & \textbf{Size} & \textbf{Purpose} \\
\hline
Preamble & 8 B & 7 bytes \texttt{10101010} $+$ 1 byte \texttt{10101011}; sync receiver clock \\
\hline
Dest.\ MAC & 6 B & If it matches (or is broadcast) pass up to the network layer, else discard \\
\hline
Source MAC & 6 B & Sending adapter's address \\
\hline
Type & 2 B & Upper-layer protocol (mostly IP) --- for demux at the receiver \\
\hline
Data (payload) & 46--1500 B & The IP datagram \\
\hline
CRC & 4 B & Error detection; frame dropped if it fails \\
\hline
\end{tabularx}

\medskip
\noindent Ethernet is \textbf{connectionless} (no handshake between NICs) and
\textbf{unreliable} (the receiving NIC sends no ACK/NAK --- data in dropped
frames is recovered only if a higher layer such as TCP retransmits). Its MAC
protocol is \textbf{unslotted CSMA/CD with binary backoff}. The \textbf{802.3}
family shares one MAC protocol and frame format across many physical
media/speeds (100BASE-TX copper, 100BASE-FX fiber, \dots).

\subsection{Switches}

A \textbf{switch} is an active link-layer device: it \emph{stores and
forwards} Ethernet frames, examining each incoming frame's MAC address to
selectively forward it to one or more outgoing links. It is \textbf{transparent}
(hosts are unaware of it) and \textbf{plug-and-play / self-learning} (no
configuration). Because each host has a dedicated full-duplex link to the
switch, there are no collisions and \textbf{multiple transmissions can happen
simultaneously} --- e.g.\ A$\to$A$'$ and B$\to$B$'$ at once (but not two frames
to the same output A$'$). Each link is its own collision domain.

\subsubsection{Switch table and self-learning}
The switch keeps a \textbf{switch table} of
$\langle$MAC address, interface, timestamp$\rangle$ (it resembles a routing
table but is built differently). It \emph{learns} by recording the
\emph{source} MAC and arriving interface of each frame. Forwarding logic:

\begin{lstlisting}[caption={Switch frame filtering / forwarding}]
When a frame arrives on an interface:
  1. record (source MAC, arriving interface) in the switch table
  2. look up the destination MAC in the switch table
  3. if an entry is found:
        if it maps to the arriving interface:   drop (filter)
        else:                                    forward on that one interface
     else:
        flood  (forward on all interfaces except the arriving one)
\end{lstlisting}
Interconnected switches self-learn in exactly the same way, so a frame from $A$
to a distant $G$ is forwarded switch-to-switch toward $G$ once the relevant
locations have been learned.

\subsubsection{Switches vs.\ routers}
Both are store-and-forward. A \textbf{switch} is a Layer-2 device: it examines
link-layer headers and builds its table by \emph{flooding + learning} on MAC
addresses. A \textbf{router} is a Layer-3 device: it examines network-layer
headers and computes tables with \emph{routing algorithms} on IP addresses,
and it isolates broadcast domains.

\subsection{VLANs}

\textbf{Motivation.} A single switched LAN is one broadcast domain: all Layer-2
broadcast traffic (ARP, DHCP, unknown-MAC floods) must cross the entire LAN
(scaling, security, and privacy issues), and moving a user's office can force
an awkward physical/logical mismatch (a CS user physically on the EE switch but
wanting to stay logically on CS).

\textbf{Port-based VLANs.} A VLAN-capable switch is configured so its ports are
grouped into isolated virtual LANs --- one physical switch behaves as several
virtual switches (e.g.\ ports 1--8 $=$ EE, ports 9--15 $=$ CS).
\textbf{Traffic isolation:} frames to/from ports 1--8 reach only ports 1--8.
\textbf{Dynamic membership:} ports can be reassigned among VLANs.
\textbf{Inter-VLAN traffic} goes through a router (just as with separate
switches; vendors sell combined switch+router boxes).

\textbf{VLANs across switches --- trunking.} A \textbf{trunk port} carries
frames of \emph{multiple} VLANs between switches. Such frames need a VLAN ID,
added by \textbf{802.1Q}: a 4-byte tag is inserted after the source-MAC field
--- a 2-byte Tag Protocol Identifier (\texttt{0x8100}) plus a 2-byte Tag
Control field (a 12-bit VLAN ID and a 3-bit priority, like IP ToS) --- and the
CRC is recomputed.

\subsection{Link Virtualization and Data Centers (Brief)}

\textbf{MPLS} (Multi-Protocol Label Switching) forwards on short fixed-length
\emph{labels} rather than IP longest-prefix match, enabling faster forwarding
and traffic engineering. \textbf{Data-center networks} interconnect tens of
thousands of hosts using hierarchical topologies (access / aggregation / core
tiers, often with rich multipath) for high bisection bandwidth.

\subsection{A Day in the Life of a Web Request (Synthesis)}

Fetching \texttt{www.google.com} exercises the whole stack at once:
\begin{enumerate}
    \item \textbf{DHCP} (UDP/IP/Ethernet, broadcast): the laptop obtains its IP
    address, subnet mask, first-hop router, and DNS server.
    \item \textbf{ARP:} to reach the DNS server / first-hop router, the host
    ARPs for the router's MAC.
    \item \textbf{DNS} (UDP): resolve the name (root $\to$ TLD $\to$
    authoritative, via the local resolver).
    \item \textbf{TCP} 3-way handshake to the web server (port 80/443).
    \item \textbf{HTTP} GET; the response returns down the stack
    (HTTP$\to$TCP$\to$IP$\to$Ethernet), through switches (self-learning,
    Layer-2) and routers (LPM forwarding, Layer-2 rewriting) along the way.
\end{enumerate}

% ==================================================================
\section{Cheat Sheet}
% ==================================================================

\subsection{Key Formulas}

\renewcommand{\arraystretch}{1.45}
\noindent\begin{tabularx}{\textwidth}{|l|X|}
\hline
\textbf{Quantity} & \textbf{Formula} \\
\hline
Transmission delay & $d_{\text{trans}} = L/R$ \quad ($L$ bits, rate $R$ bps) \\
\hline
Propagation delay & $d_{\text{prop}} = d/s$ \quad ($s \approx 2\times 10^{8}$ m/s) \\
\hline
Nodal delay & $d_{\text{nodal}} = d_{\text{proc}} + d_{\text{queue}} + d_{\text{trans}} + d_{\text{prop}}$ \\
\hline
Traffic intensity & $La/R$ \; ($<1$ good, $\to 1$ large delay, $>1$ unbounded) \\
\hline
End-to-end throughput & $\min(R_s, R_c)$; shared: $\min(R_s, R_c, R/N)$ \\
\hline
Packet-switch capacity (binomial) & $P(X{=}k)=\binom{n}{k}p^{k}(1-p)^{n-k}$ \\
\hline
Non-persistent HTTP (per object) & $2\,RTT + T_{\text{transmit}}$ \\
\hline
File dist.\ (client--server) & $D_{\text{c-s}} \geq \max\{NF/u_s,\; F/d_{\min}\}$ \\
\hline
File dist.\ (P2P) & $D_{\text{P2P}} \geq \max\{F/u_s,\; F/d_{\min},\; NF/(u_s{+}\textstyle\sum u_i)\}$ \\
\hline
EWMA RTT & $\text{EstRTT}=(1{-}\alpha)\text{EstRTT}+\alpha\,\text{SampleRTT}$, $\alpha{=}0.125$ \\
\hline
RTT deviation & $\text{DevRTT}=(1{-}\beta)\text{DevRTT}+\beta|\text{SampleRTT}-\text{EstRTT}|$, $\beta{=}0.25$ \\
\hline
TCP timeout & $\text{TimeoutInterval}=\text{EstRTT}+4\,\text{DevRTT}$ \\
\hline
Stop-and-wait utilization & $U = \dfrac{L/R}{RTT + L/R}$ \\
\hline
Pipelined ($n$ packets) & $U = \dfrac{n\,L/R}{RTT + L/R}$ \\
\hline
TCP rate & $\approx \texttt{cwnd}/RTT$ \\
\hline
TCP avg.\ throughput & $\tfrac{3}{4}\,W/RTT$ \\
\hline
TCP throughput vs.\ loss (Mathis) & $1.22\cdot\text{MSS}\,/\,(RTT\sqrt{L})$ \\
\hline
SR window constraint & window size $\leq \tfrac{1}{2}\times$ seq-number space \\
\hline
Router switching rate (ideal) & $N \times R$ \; ($N$ ports, line rate $R$) \\
\hline
Router buffer (RFC 3439) & $\text{RTT} \times C$ \\
\hline
Router buffer ($N$ flows) & $\text{RTT} \times C / \sqrt{N}$ \\
\hline
WFQ per-class share & $w_i / \textstyle\sum_j w_j$ \\
\hline
Hosts in a \texttt{/x} subnet & $2^{\,32-x} - 2$ \; (minus network + broadcast) \\
\hline
Fragment offset & (bytes before this fragment) $/\,8$ \\
\hline
Max frag.\ payload & $\lfloor (\text{MTU} - 20)/8 \rfloor \times 8$ bytes \\
\hline
Bellman-Ford eqn.\ & $d_x(y) = \min_v\{\,c(x,v) + d_v(y)\,\}$ \\
\hline
Dijkstra complexity & $O(n^2)$, or $O(n\log n)$ with a heap \\
\hline
Slotted ALOHA (max eff.) & $1/e \approx 0.37$ \\
\hline
Pure ALOHA (max eff.) & $1/(2e) \approx 0.18$ \\
\hline
CSMA/CD efficiency & $\approx \dfrac{1}{1 + 5\,t_{\text{prop}}/t_{\text{trans}}}$ \\
\hline
CRC remainder & $R = \text{remainder}(D\!\cdot\!2^{r} / G)\ \bmod 2$ \; ($r{+}1$-bit $G$) \\
\hline
\end{tabularx}

\subsection{Well-Known Ports}

\noindent\begin{tabularx}{\textwidth}{|l|l|X|}
\hline
\textbf{Port} & \textbf{Protocol} & \textbf{Transport / notes} \\
\hline
20/21 & FTP (data/control) & TCP \\
\hline
22 & SSH & TCP \\
\hline
25 & SMTP & TCP (mail \emph{between} servers) \\
\hline
53 & DNS & UDP (TCP for large transfers) \\
\hline
80 & HTTP & TCP \\
\hline
143 & IMAP & TCP (mail retrieval) \\
\hline
161 & SNMP & UDP \\
\hline
443 & HTTPS / HTTP/3 & TCP (TLS) / UDP (QUIC) \\
\hline
\end{tabularx}

\subsection{HTTP Methods and Status Codes}

\noindent\begin{tabularx}{\textwidth}{|l|X||l|X|}
\hline
\textbf{Method} & \textbf{Purpose} & \textbf{Code} & \textbf{Meaning} \\
\hline
GET  & Retrieve object & 200 & OK \\
\hline
POST & Send form data in body & 301 & Moved Permanently \\
\hline
HEAD & Headers only & 304 & Not Modified (conditional GET) \\
\hline
PUT  & Upload / replace object & 400 & Bad Request \\
\hline
     &  & 404 & Not Found \\
\hline
     &  & 505 & HTTP Version Not Supported \\
\hline
\end{tabularx}

\subsection{DNS Record Types}

\noindent\begin{tabularx}{\textwidth}{|l|X|}
\hline
\textbf{Type} & \textbf{Maps / does} \\
\hline
A     & name $\to$ IPv4 address \\
\hline
AAAA  & name $\to$ IPv6 address \\
\hline
NS    & delegate a domain to a name server \\
\hline
CNAME & alias $\to$ canonical name \\
\hline
MX    & domain $\to$ mail server (lower priority number tried first) \\
\hline
\end{tabularx}

\subsection{TCP vs.\ UDP}

\noindent\begin{tabularx}{\textwidth}{|l|X|X|}
\hline
 & \textbf{TCP} & \textbf{UDP} \\
\hline
Connection & Yes (3-way handshake) & No \\
\hline
Reliability & Reliable, in-order byte stream & Best effort, may lose/reorder \\
\hline
Flow control & Yes (\texttt{rwnd}) & No \\
\hline
Congestion control & Yes (AIMD / CUBIC / BBR) & No \\
\hline
Header & 20+ bytes & 8 bytes \\
\hline
Demux key & 4-tuple (src/dst IP + port) & 2-tuple (dst IP + port) \\
\hline
Typical use & Web, email, file transfer & DNS, SNMP, streaming, HTTP/3 \\
\hline
\end{tabularx}

\subsection{Circuit vs.\ Packet Switching}

\noindent\begin{tabularx}{\textwidth}{|l|X|X|}
\hline
 & \textbf{Circuit} & \textbf{Packet} \\
\hline
Resources & Reserved (dedicated) & Shared on demand \\
\hline
Setup & Required & None \\
\hline
Performance & Guaranteed & Best effort (delay/loss) \\
\hline
Efficiency & Idle circuits waste capacity & Statistical multiplexing \\
\hline
Used by & Traditional telephone net & The Internet \\
\hline
\end{tabularx}

\subsection{Go-Back-N vs.\ Selective Repeat}

\noindent\begin{tabularx}{\textwidth}{|l|X|X|}
\hline
 & \textbf{Go-Back-N} & \textbf{Selective Repeat} \\
\hline
ACK type & Cumulative & Individual \\
\hline
Receiver buffering & None (discard out-of-order) & Buffers out-of-order \\
\hline
Timers & One (oldest packet) & One per unACKed packet \\
\hline
On loss & Resend window from $n$ onward & Resend only the lost packet \\
\hline
\end{tabularx}

\subsection{The Five-Layer Stack (with units and example protocols)}

\noindent\begin{tabularx}{\textwidth}{|l|l|X|}
\hline
\textbf{Layer} & \textbf{Unit} & \textbf{Examples} \\
\hline
5. Application & message  & HTTP, SMTP, IMAP, DNS \\
\hline
4. Transport   & segment  & TCP, UDP \\
\hline
3. Network     & datagram & IP (the ``narrow waist''), routing protocols \\
\hline
2. Link        & frame    & Ethernet, 802.11 (WiFi), PPP \\
\hline
1. Physical    & bit      & copper, fiber, radio \\
\hline
\end{tabularx}

\subsection{TCP Congestion-Control State Reactions}

\noindent\begin{tabularx}{\textwidth}{|l|X|}
\hline
\textbf{Event} & \textbf{Action} \\
\hline
Connection start & \texttt{cwnd}$=1$ MSS; slow start (double per RTT) \\
\hline
\texttt{cwnd} $\geq$ \texttt{ssthresh} & Enter congestion avoidance ($+1$ MSS / RTT) \\
\hline
Triple duplicate ACK (Reno) & \texttt{ssthresh}$=$\texttt{cwnd}$/2$; \texttt{cwnd}$=$\texttt{ssthresh}$+3$ (fast recovery) \\
\hline
Timeout (Tahoe) & \texttt{ssthresh}$=$\texttt{cwnd}$/2$; \texttt{cwnd}$=1$ MSS; slow start \\
\hline
\end{tabularx}

\subsection{Router Data Plane vs.\ Control Plane}

\noindent\begin{tabularx}{\textwidth}{|l|X|X|}
\hline
 & \textbf{Data plane} & \textbf{Control plane} \\
\hline
Scope & Local, per-router & Network-wide \\
\hline
Job & Forwarding (input $\to$ output port) & Routing (build tables) \\
\hline
Implemented in & Hardware & Software \\
\hline
Timescale & Nanoseconds & Milliseconds \\
\hline
Approaches & Match+action lookup & Per-router alg.\ \emph{or} SDN controller \\
\hline
\end{tabularx}

\subsection{Packet Scheduling Policies}

\noindent\begin{tabularx}{\textwidth}{|l|X|X|}
\hline
\textbf{Policy} & \textbf{Rule} & \textbf{Property} \\
\hline
FCFS / FIFO & Serve in arrival order & Simple; no isolation \\
\hline
Priority & Highest non-empty class first (FCFS within) & Low-priority starvation risk \\
\hline
Round Robin & One packet per class, cyclically & Fair-ish, no weights \\
\hline
WFQ & Class $i$ gets $w_i/\sum_j w_j$ share & Per-class min-bandwidth guarantee \\
\hline
\end{tabularx}

\subsection{IPv4 Header Fields (Quick Reference)}

\noindent\begin{tabularx}{\textwidth}{|l|X|}
\hline
\textbf{Field} & \textbf{Purpose} \\
\hline
Version / Header len & IP version; header size in 32-bit words \\
\hline
Type of Service & Diffserv (bits 0--5) + ECN (bits 6--7) \\
\hline
Total length & Datagram size in bytes (max 64 KB) \\
\hline
Identifier / Flags / Offset & Fragmentation and reassembly \\
\hline
TTL & Max remaining hops (decremented per router) \\
\hline
Upper-layer protocol & Demux key (6 = TCP, 17 = UDP) \\
\hline
Header checksum & Error check over the header only \\
\hline
Source / Dest address & 32-bit endpoints \\
\hline
\end{tabularx}

\subsection{CIDR / Subnet Mask Quick Reference}

\noindent\begin{tabularx}{\textwidth}{|l|l|l|X|}
\hline
\textbf{Prefix} & \textbf{Mask} & \textbf{Addresses} & \textbf{Usable hosts ($2^{32-x}{-}2$)} \\
\hline
/8  & 255.0.0.0       & $2^{24}$ & 16{,}777{,}214 \\
\hline
/16 & 255.255.0.0     & $2^{16}$ & 65{,}534 \\
\hline
/20 & 255.255.240.0   & $2^{12}$ & 4{,}094 \\
\hline
/23 & 255.255.254.0   & $2^{9}$  & 510 \\
\hline
/24 & 255.255.255.0   & $2^{8}$  & 254 \\
\hline
/30 & 255.255.255.252 & $2^{2}$  & 2 (point-to-point links) \\
\hline
\end{tabularx}

\subsection{Private and Special IPv4 Addresses}

\noindent\begin{tabularx}{\textwidth}{|l|X|}
\hline
\textbf{Block / address} & \textbf{Use} \\
\hline
\texttt{10.0.0.0/8} & Private (RFC 1918) \\
\hline
\texttt{172.16.0.0/12} & Private (RFC 1918) \\
\hline
\texttt{192.168.0.0/16} & Private (RFC 1918) \\
\hline
\texttt{127.0.0.0/8} & Loopback \\
\hline
\texttt{0.0.0.0} & ``This host'' / unspecified (DHCP source) \\
\hline
\texttt{255.255.255.255} & Limited broadcast (DHCP destination) \\
\hline
Ports 67 / 68 & DHCP server / client (over UDP) \\
\hline
\end{tabularx}

\subsection{IPv4 vs.\ IPv6}

\noindent\begin{tabularx}{\textwidth}{|l|X|X|}
\hline
 & \textbf{IPv4} & \textbf{IPv6} \\
\hline
Address size & 32-bit & 128-bit \\
\hline
Header & Variable ($\geq$ 20 B) & Fixed 40 B \\
\hline
Checksum & Yes & No (speeds forwarding) \\
\hline
Fragmentation & Routers \emph{and} host & Host only (not in-network) \\
\hline
Options & In header & ``Next-header'' chain \\
\hline
Notation & Dotted decimal & Colon-hex \\
\hline
Address setup & DHCP / static & SLAAC / DHCPv6 \\
\hline
\end{tabularx}

\subsection{DHCP (DORA) and NAT}

\noindent\begin{tabularx}{\textwidth}{|l|X|}
\hline
\textbf{Step / concept} & \textbf{Meaning} \\
\hline
Discover & Client broadcasts to find a server \\
\hline
Offer & Server proposes address (\texttt{yiaddr}) + lease \\
\hline
Request & Client asks for the offered address \\
\hline
ACK & Server confirms; also returns router, DNS, mask \\
\hline
NAT out & (src IP, port) $\to$ (public IP, new port), record in table \\
\hline
NAT in & (public IP, new port) $\to$ stored (src IP, port) \\
\hline
\end{tabularx}

\subsection{Routing Algorithms and Protocols}

\noindent\begin{tabularx}{\textwidth}{|l|X|X|}
\hline
 & \textbf{Link-State} & \textbf{Distance-Vector} \\
\hline
Knows & Full topology & Neighbors + their vectors \\
\hline
Computes with & Dijkstra & Bellman-Ford \\
\hline
Talks to & All (flooding) & Neighbors only \\
\hline
Weakness & Oscillations & Count-to-infinity \\
\hline
Example & OSPF & RIP \\
\hline
\end{tabularx}

\medskip
\noindent\begin{tabularx}{\textwidth}{|l|X|X|}
\hline
 & \textbf{Intra-AS (IGP)} & \textbf{Inter-AS (BGP)} \\
\hline
Scope & Within one AS & Between ASes \\
\hline
Type & OSPF (LS), RIP (DV) & Path-vector \\
\hline
Goal & Least cost & Policy (economics) \\
\hline
Key attrs & Link costs, areas & AS-PATH, NEXT-HOP, local-pref \\
\hline
Selection & Shortest path & Local-pref $>$ short AS-PATH $>$ hot-potato \\
\hline
\end{tabularx}

\subsection{ICMP Common Messages}

\noindent\begin{tabularx}{\textwidth}{|l|l|X|}
\hline
\textbf{Type} & \textbf{Name} & \textbf{Used by / meaning} \\
\hline
8 / 0 & Echo request/reply & \texttt{ping} \\
\hline
3 & Destination unreachable & network/host/port/frag-needed (codes) \\
\hline
11 & TTL expired & \texttt{traceroute} \\
\hline
\end{tabularx}

\subsection{Multiple Access Protocols}

\noindent\begin{tabularx}{\textwidth}{|l|X|X|}
\hline
\textbf{Family} & \textbf{Protocols} & \textbf{Idea} \\
\hline
Partitioning & TDMA, FDMA & Fixed slots/bands; no collisions, low-load waste \\
\hline
Random access & ALOHA, CSMA/CD, CSMA/CA & Contend; detect/avoid \& recover from collisions \\
\hline
Taking turns & Polling, token passing & Coordinate whose turn it is \\
\hline
\end{tabularx}

\subsection{Error Detection Methods}

\noindent\begin{tabularx}{\textwidth}{|l|X|}
\hline
\textbf{Method} & \textbf{Strength} \\
\hline
1-bit parity & Detects odd \# of bit errors \\
\hline
2-D parity & Detects + \emph{corrects} a single-bit error \\
\hline
Internet checksum & Weak; software (transport layer) \\
\hline
CRC & Strong; hardware; detects bursts $\leq r$ bits (Ethernet, WiFi) \\
\hline
\end{tabularx}

\subsection{Link Layer: MAC vs.\ IP, Switch vs.\ Router}

\noindent\begin{tabularx}{\textwidth}{|l|X|X|}
\hline
 & \textbf{MAC address} & \textbf{IP address} \\
\hline
Size & 48-bit & 32-bit (v4) \\
\hline
Structure & Flat & Hierarchical \\
\hline
Scope & One LAN & End-to-end \\
\hline
Assignment & Burned into NIC (portable) & Subnet-dependent (DHCP/static) \\
\hline
Resolved by & --- & ARP (IP $\to$ MAC on a LAN) \\
\hline
\end{tabularx}

\medskip
\noindent\begin{tabularx}{\textwidth}{|l|X|X|}
\hline
 & \textbf{Switch (L2)} & \textbf{Router (L3)} \\
\hline
Uses & MAC addresses & IP addresses \\
\hline
Table built by & Self-learning (source MAC) & Routing algorithms \\
\hline
Unknown dest & Flood & Drop / default route \\
\hline
Broadcast domain & Does \emph{not} isolate & Isolates \\
\hline
Config & Plug-and-play & Needs addressing/routing \\
\hline
\end{tabularx}

\subsection{Ethernet Frame Fields}

\noindent\begin{tabularx}{\textwidth}{|l|l|X|}
\hline
\textbf{Field} & \textbf{Size} & \textbf{Purpose} \\
\hline
Preamble & 8 B & Clock synchronization \\
\hline
Dest / Source MAC & 6 B each & Link-layer addresses \\
\hline
Type & 2 B & Upper-layer demux (usually IP) \\
\hline
Data (payload) & 46--1500 B & The IP datagram \\
\hline
CRC & 4 B & Error detection (frame dropped on failure) \\
\hline
\end{tabularx}

\subsection{Intra-AS Routing Protocols}

\noindent\begin{tabularx}{\textwidth}{|l|X|X|}
\hline
\textbf{Protocol} & \textbf{Type} & \textbf{Notes} \\
\hline
RIP & Distance-vector & DVs every 30 s; mostly historical \\
\hline
EIGRP & Distance-vector & Cisco, opened 2013 \\
\hline
OSPF / IS-IS & Link-state (Dijkstra) & Floods LSAs over IP; authenticated; areas \\
\hline
\end{tabularx}

\subsection{BGP: Attributes and Route Selection}

\noindent\begin{tabularx}{\textwidth}{|l|X|}
\hline
\textbf{Concept} & \textbf{Meaning} \\
\hline
eBGP / iBGP & Learn reachability from other ASes / propagate it inside the AS \\
\hline
Transport & Semi-permanent TCP session between peers \\
\hline
AS-PATH & List of ASes the advert crossed (loop detection) \\
\hline
NEXT-HOP & Internal router toward the next-hop AS \\
\hline
Messages & OPEN, UPDATE, KEEPALIVE, NOTIFICATION \\
\hline
\end{tabularx}

\medskip
\noindent\textbf{Route selection order:} (1) highest local-preference (policy);
(2) shortest AS-PATH; (3) closest NEXT-HOP (hot-potato); (4) other tie-breakers.
\textbf{Hot-potato:} pick the gateway with least \emph{intra}-domain cost.

\subsection{ARP and Switch Self-Learning}

\noindent\begin{tabularx}{\textwidth}{|l|X|}
\hline
\textbf{Step / rule} & \textbf{Action} \\
\hline
ARP query & Broadcast (dest MAC \texttt{FF-FF-FF-FF-FF-FF}) with target IP \\
\hline
ARP reply & Owner unicasts its MAC; requester caches $\langle$IP,MAC,TTL$\rangle$ ($\sim$20 min) \\
\hline
Switch learn & Record (source MAC, arriving interface) \\
\hline
Dest known, same port & Drop (filter) \\
\hline
Dest known, other port & Forward on that one interface \\
\hline
Dest unknown / broadcast & Flood (all ports except arriving) \\
\hline
\end{tabularx}

\subsection{VLANs and 802.1Q}

\noindent\begin{tabularx}{\textwidth}{|l|X|}
\hline
\textbf{Concept} & \textbf{Meaning} \\
\hline
Port-based VLAN & One physical switch acts as several virtual switches \\
\hline
Traffic isolation & Frames stay within their VLAN's ports \\
\hline
Inter-VLAN & Requires a router (Layer 3) \\
\hline
Trunk port & Carries multiple VLANs between switches \\
\hline
802.1Q tag & 4 B: TPID \texttt{0x8100} + 12-bit VLAN ID + 3-bit priority; CRC recomputed \\
\hline
\end{tabularx}

\vfill
\noindent\textit{\color{Thistle}These notes paraphrase and reorganize the
lecture material (Kurose \& Ross, 8th ed.; UC Berkeley CS\,168) into a single
study reference spanning the full exam scope --- Transport (Ch.\,3),
Network-layer Data Plane (Ch.\,4) and Control Plane (Ch.\,5), and the Link
Layer (Ch.\,6) --- with added code (sockets, SMTP, checksum, fragmentation,
subnets, Dijkstra, distance-vector, CRC, CSMA/CD, switch learning) and an
expanded cheat sheet.}

\end{document}
